\section{AIME I 2020}
        \begin{enumerate}
        	\item (\textit{AIME I 2020}) In $\triangle ABC$ with $AB=AC,$ point $D$ lies strictly between $A$ and $C$ on side $\overline{AC},$ and point $E$ lies strictly between $A$ and $B$ on side $\overline{AB}$ such that $AE=ED=DB=BC.$ The degree measure of $\angle ABC$ is $\frac{m}{n},$ where $m$ and $n$ are relatively prime positive integers. Find $m+n.$


 
 \textbf{Answer: }547



\textbf{Solution 1}


                    \begin{figure}[H]
                        \centering
                        \includegraphics[width=0.5\textwidth]{Solution/AIME_2020_I_1_Solution_Figure_1.png}
                    \end{figure}
                    

 If we set $\angle{BAC}$ to $x$, we can find all other angles through these two properties:1. Angles in a triangle sum to $180^{\circ}$.2. The base angles of an isosceles triangle are congruent.

 Now we angle chase. $\angle{ADE}=\angle{EAD}=x$, $\angle{AED} = 180-2x$, $\angle{BED}=\angle{EBD}=2x$, $\angle{EDB} = 180-4x$, $\angle{BDC} = \angle{BCD} = 3x$, $\angle{CBD} = 180-6x$. Since $AB = AC$ as given by the problem, $\angle{ABC} = \angle{ACB}$, so $180-4x=3x$. Therefore, $x = 180/7^{\circ}$, and our desired angle is 
 \[180-4\left(\frac{180}{7}\right) = \frac{540}{7}\] for an answer of $\boxed{547}$.

 





\textbf{Solution 2}

Let $\angle{BAC}$ be $x$ in degrees. $\angle{ADE}=x$.By Exterior Angle Theorem on triangle $AED$, $\angle{BED}=2x$.By Exterior Angle Theorem on triangle $ADB$, $\angle{BDC}=3x$.This tells us $\angle{BCA}=\angle{ABC}=3x$ and $3x+3x+x=180$.Thus $x=\frac{180}{7}$ and we want $\angle{ABC}=3x=\frac{540}{7}$ to get an answer of $\boxed{547}$.

 





\textbf{Solution 3 (Official MAA)}

Let $x = \angle ABC = \angle ACB$. Because $\triangle BCD$ is isosceles, $\angle CBD = 180^\circ - 2x$. Then
 \[\angle DBE = x - \angle CBD = x - (180^\circ - 2x) =  3x - 180^\circ\!.\]Because $\triangle EDA$ and $\triangle DBE$ are also isosceles,
 \[\angle BAC =\frac12(\angle EAD + \angle ADE) = \frac12(\angle BED)= \frac12(\angle DBE)\]
 \[= \frac12 (3x - 180^\circ) = \frac32x-90^\circ\!.\]Because $\triangle ABC$ is isosceles, $\angle BAC$ is also $180^\circ-2x$, so $\frac32x - 90^\circ = 180^\circ - 2x$, and it follows that$\angle ABC = x = \left(\frac{540}7\right)^\circ$. The requested sum is $540+7 = 547$.

 
                    \begin{figure}[H]
                        \centering
                        \includegraphics[width=0.5\textwidth]{Solution/AIME_2020_I_1_Solution_Figure_1.png}
                    \end{figure}
                    

 \href{https://artofproblemsolving.com/wiki/index.php/1961\_AHSME\_Problems/Problem\_25}{https://artofproblemsolving.com/wiki/index.php/1961\_AHSME\_Problems/Problem\_25}(Almost Mirrored)

 See here for a video solution:

 





\textbf{Solution 4 (writing equations)}

graph soon

 We write equations based on the triangle sum of angles theorem. There are angles that do not need variables as the less variables the better.

 
 \begin{align*} \angle A &= y \\ \angle B &= x+z \\ \angle C &= \frac{180-x}{2} \\ \end{align*}


 Then, using triangle sum of angles theorem, we find that

 
 \begin{align*} \angle A + \angle B + \angle C = x+y+z+\frac{180-x}{2}=180 \\ \end{align*}


 Now we just need to find the variables.

 
 \begin{align*} (180-2y)+z = 180& \\ (180-2z)+y+\frac{180-x}{2} = 180& \\ \end{align*}


 Notice how all the equations equal 180. We can use this to write

 
 \begin{align*} (180-2y)+z = (180-2z)+y+\frac{180-x}{2}=x+y+z+\frac{180-x}{2} \\ \end{align*}


 Simplifying, we get

 
 \begin{align*} (180-2y)+z=(180-2z)+y+\frac{180-x}{2} \\ 360-4y+2z=360-4z+2y+180-x \\ \end{align*}


 
 \begin{align*} 6z=6y+180-x \\ x=6y-6z+180 \\ \end{align*}


 
 \begin{align*} (180-2y)+z=6y-6z+180+y+z+\frac{180-(6y-6+180)}{2} \\ 360-4y+2z=12y-12z+360+2y+2z+180-6y+6z-180 \\ \end{align*}


 
 \begin{align*} 6z=12y& \\ z=2y& \\ \end{align*}


 Theres more. We are at a dead end right now because we forgot that the problem states that the triangle is isosceles. With this, we can write the equation

 
 \begin{align*} \frac{180-x}{2}=x+z \\ \end{align*}


 Substituting $z$ with $2y$, we get

 
 \begin{align*} \frac{180-x}{2}=x+2y \\ 180-x=2x+4y \\ \end{align*}

 \begin{align*} 180-(6y-6z+180)=2(6y-6z+180)+4y& \\ 180-6y+12y-180=12y-24y+360+4y& \\ \end{align*}

 \begin{align*} 6y=-8y+360& \\ \end{align*}


 With this, we get

 
 \begin{align*} y=\frac{180}{7} \\ x=\frac{180}{7} \\ z=\frac{360}{7} \\ \end{align*}


 And a final answer of $\frac{180}{7}+\frac{360}{7} = \frac{540}{7} = \boxed{547}$.

 





\textbf{Solution without words}


                    \begin{figure}[H]
                        \centering
                        \includegraphics[width=0.5\textwidth]{Solution/AIME_2020_I_1_Solution_Figure_1.png}
                    \end{figure}
                    
                    \begin{figure}[H]
                        \centering
                        \includegraphics[width=0.5\textwidth]{Solution/AIME_2020_I_1_Solution_Figure_1.png}
                    \end{figure}
                    \textbf{vladimir.shelomovskii@gmail.com, vvsss}

 




	\item (\textit{AIME I 2020}) There is a unique positive real number $x$ such that the three numbers $\log_8(2x),\log_4x,$ and $\log_2x,$ in that order, form a geometric progression with positive common ratio. The number $x$ can be written as $\frac{m}{n},$ where $m$ and $n$ are relatively prime positive integers. Find $m+n.$


 
 \textbf{Answer: }017



\textbf{Solution 1}

Since these form a geometric series, $\frac{\log_2{x}}{\log_4{x}}$ is the common ratio. Rewriting this, we get $\frac{\log_x{4}}{\log_x{2}} = \log_2{4} = 2$ by base change formula. Therefore, the common ratio is 2. Now $\frac{\log_4{x}}{\log_8{2x}} = 2 \implies \log_4{x} = 2\log_8{2} + 2\log_8{x} \implies \frac{1}{2}\log_2{x} = \frac{2}{3} + \frac{2}{3}\log_2{x}$

 $\implies -\frac{1}{6}\log_2{x} = \frac{2}{3} \implies \log_2{x} = -4 \implies x = \frac{1}{16}$. Therefore, $1 + 16 = \boxed{017}$.

 





\textbf{Solution 2}

If we set $x=2^y$, we can obtain three terms of a geometric sequence through logarithm properties. The three terms are 
 \[\frac{y+1}{3}, \frac{y}{2}, y.\] In a three-term geometric sequence, the middle term squared is equal to the product of the other two terms, so we obtain the following: 
 \[\frac{y^2+y}{3} = \frac{y^2}{4},\] which can be solved to reveal $y = -4$. Therefore, $x = 2^{-4} = \frac{1}{16}$, so our answer is $\boxed{017}$.

 -molocyxu

 





\textbf{Solution 3}

Let $r$ be the common ratio. We have 
 \[r = \frac{\log_4{(x)}}{\log_8{(2x)}} = \frac{\log_2{(x)}}{\log_4{(x)}}\]Hence we obtain 
 \[(\log_4{(x)})(\log_4{(x)}) = (\log_8{(2x)})(\log_2{(x)})\]Ideally we change everything to base $64$ and we can get: 
 \[(\log_{64}{(x^3)})(\log_{64}{(x^3)}) = (\log_{64}{(x^6)})(\log_{64}{(4x^2)})\]Now divide to get: 
 \[\frac{\log_{64}{(x^3)}}{\log_{64}{(4x^2)}} = \frac{\log_{64}{(x^6)}}{\log_{64}{(x^3)}}\]By change-of-base we obtain: 
 \[\log_{(4x^2)}{(x^3)} = \log_{(x^3)}{(x^6)} = 2\]Hence $(4x^2)^2 = x^3 \rightarrow 16x^4 = x^3 \rightarrow x = \frac{1}{16}$ and we have $1+16 = \boxed{017}$ as desired.

 





\textbf{Solution 4 (Exponents > Logarithms)}

Let $r$ be the common ratio, and let $a$ be the starting term ($a=\log_{8}{(2x)}$). We then have: 
 \[\log_{8}{(2x)}=a, \log_{4}{(x)}=ar, \log_{2}{(x)}=ar^2\] Rearranging these equations gives: 
 \[8^a=2x, 4^{ar}=x, 2^{ar^2}=x\]Deal with the last two equations first: Setting them equal gives: 
 \[4^{ar}=2^{ar^2} \implies 2^{2ar}=2^{ar^2} \implies 2ar=ar^2 \implies r=2\] Using this value of $r$, substitute into the first and second equations (or the first and third, it doesn't really matter) to get: 
 \[8^a=2x, 4^{2a}=x\] Changing these to a common base gives: 
 \[2^{3a}=2x, 2^{4a}=x\] Dividing the first equation by 2 on both sides yields: 
 \[2^{3a-1}=x\] Setting these equations equal to each other and removing the exponent again gives: 
 \[3a-1=4a \implies a=-1\] Substituting this back into the first equation gives: 
 \[8^{-1}=2x \implies 2x=\frac{1}{8} \implies x=\frac{1}{16}\] Therefore, $m+n=1+16=\boxed{017}$

 





\textbf{Solution 5}

We can relate the logarithms as follows:

 
 \[\frac{\log_4{x}}{\log_8{(2x)}}=\frac{\log_2{x}}{\log_4{x}}\]
 \[\log_8{(2x)}\log_2{x}=\log_4{x}\log_4{x}\]

 Now we can convert all logarithm bases to $2$ using the identity $\log_a{b}=\log_{a^c}{b^c}$:

 
 \[\log_2{\sqrt[3]{2x}}\log_2{x}=\log_2{\sqrt{x}}\log_2{\sqrt{x}}\]

 We can solve for $x$ as follows:

 
 \[\frac{1}{3}\log_2{(2x)}\log_2{x}=\frac{1}{4}\log_2{x}\log_2{x}\]
 \[\frac{1}{3}\log_2{(2x)}=\frac{1}{4}\log_2{x}\]
 \[\frac{1}{3}\log_2{2}+\frac{1}{3}\log_2{x}=\frac{1}{4}\log_2{x}\]We get $x=\frac{1}{16}$. Verifying that the common ratio is positive, we find the answer of $\boxed{017}$.

 


 





\textbf{Solution 6}

If the numbers are in a geometric sequence, the middle term must be the geometric mean of the surrounding terms. We can rewrite the first two logarithmic expressions as $\frac{1+\log_2{x}}{3}$ and $\frac{1}{2}\log_2{x}$, respectively. Therefore:
 \[\frac{1}{2}\log_2{x}=\sqrt{\left(\frac{1+\log_2{x}}{3}\right)\left(\log_2{x}\right)}\]Let $n=\log_2{x}$. We can rewrite the expression as:
 \[\frac{n}{2}=\sqrt{\frac{n(n+1)}{3}}\]
 \[\frac{n^2}{4}=\frac{n(n+1)}{3}\]
 \[4n(n+1)=3n^2\]
 \[4n^2+4n=3n^2\]
 \[n^2+4n=0\]
 \[n(n+4)=0\]
 \[n=0 \text{ and } -4\]Zero does not work in this case, so we consider $n=-4$: $\log_2{x}=-4 \rightarrow x=\frac{1}{16}$. Therefore, $1+16=\boxed{017}$.

 





\textbf{Solution 7}

Again, by the Change of Base Formula, obtain that the common ratio is 2. If we let $y$ be the exponent of $\log_8 (2x)$, then we have $8^y=2x;\:4^{2y}=x;\:2^{4y}=x.$ Wee can then divide the first equation by two to have the right side be $x$. Also, $2^{4y}=\left(2^{4}\right)^y=16^y$. Setting this equal to $\frac{8^y}{2}$, we can divide the two equations to get $2^y=\frac12$. Therefore, $y=-1$. After that, we can raise $16$ to the $-1$th power to get $x=16^{-1}\Rightarrow x=\frac1{16}$. We then get our sum of $1+16=\boxed{\textbf{017}}$.

 





\textbf{Solution 8 (Official MAA)}

By the Change of Base Formula the common ratio of the progression is
 \[\frac{\log_2 x}{\log_4 x} = \frac{\hphantom{m}\log_2x\hphantom{m}}{\frac{\log_2x}{\log_24}} = 2.\]Hence $x$ must satisfy
 \[2=\frac{\log_4 x}{\log_8 (2x)}= \frac{\log_2 x}{\log_2 4} \div \frac{\log_2(2x)}{\log_28} = \frac 32\cdot \frac{\log_2x}{1+\log_2x}.\]This is equivalent to $4 + 4\log_2x = 3\log_2x$. Hence $\log_2x = -4$ and $x = \frac{1}{16}$. The requested sum is $1+16 = 17$.

 





\textbf{Solution 9}

Let the series be $a,ar,ar^2$. So, converting these logarithms into exponents we have $8^a = 2x$, $4^{ar} = x$, and $2^{ar^2} = x$. So, setting the 2nd and 3rd equations equal, we have:
 \[4^{ar} = 2^{ar^2} \Longrightarrow 2^{2ar} = 2^{ar^2} \Longrightarrow 2ar = ar^2 \Longrightarrow r = 2\]. Plugging that into our 3rd equation, we have $2^{4a} = x$. We recall that $8^a = 2x \Longrightarrow 2^{3a} = 2x$. So, since $2^{4a} = x$, $2x = 2^{4a+1}$. Now, setting $2^{3a}$, and $2^{4a+1}$ equal, we get $a = -1$, so plugging into our first equation, we have
 \[\dfrac{1}{8} = 2x \Longrightarrow x = \dfrac{1}{16} \Longrightarrow \boxed{017}\]

 -jb2015007

 




	\item (\textit{AIME I 2020}) A positive integer $N$ has base-eleven representation $\underline{a}\kern 0.1em\underline{b}\kern 0.1em\underline{c}$ and base-eight representation $\underline1\kern 0.1em\underline{b}\kern 0.1em\underline{c}\kern 0.1em\underline{a},$ where $a,b,$ and $c$ represent (not necessarily distinct) digits. Find the least such $N$ expressed in base ten.


 
 \textbf{Answer: }621



\textbf{Solution 1}

From the given information, $121a+11b+c=512+64b+8c+a \implies 120a=512+53b+7c$. Since $a$, $b$, and $c$ have to be positive, $a \geq 5$. Since we need to minimize the value of $n$, we want to minimize $a$, so we have $a = 5$. Then we know $88=53b+7c$, and we can see the only solution is $b=1$, $c=5$. Finally, $515_{11} = 621_{10}$, so our answer is $\boxed{621}$.

 





\textbf{Solution 2 (Official MAA)}

The conditions of the problem imply that $121a + 11b + c = 512 + 64b + 8 c + a$, so $120 a = 512+ 53b+7c$. The maximum digit in base eight is $7,$ and because $120a \ge 512$, it must be that $a$ is $5, 6,$ or $7.$ When $a = 5$, it follows that $600=512 + 53b+7c$, which implies that $88 = 53b+7c$. Then $b$ must be $0$ or $1.$ If $b = 0$, then $c$ is not an integer, and if $b = 1$, then $7c = 35$, so $c = 5$. Thus $N = 515_{11}$, and $N=5\cdot 121 + 1\cdot 11 + 5 = 621$. The number $637_{11} =1376_{8} = 766$ also satisfies the conditions of the problem, but $621$ is the least such number.

 Minor edits by TryhardMathlete

 




	\item (\textit{AIME I 2020}) Let $S$ be the set of positive integers $N$ with the property that the last four digits of $N$ are $2020,$ and when the last four digits are removed, the result is a divisor of $N.$ For example, $42{,}020$ is in $S$ because $4$ is a divisor of $42{,}020.$ Find the sum of all the digits of all the numbers in $S.$ For example, the number $42{,}020$ contributes $4+2+0+2+0=8$ to this total.


 
 \textbf{Answer: }093



\textbf{Solution 1}

We note that any number in $S$ can be expressed as $a(10,000) + 2,020$ for some integer $a$. The problem requires that $a$ divides this number, and since we know $a$ divides $a(10,000)$, we need that $a$ divides 2020. Each number contributes the sum of the digits of $a$, as well as $2 + 0 + 2 +0 = 4$. Since $2020$ can be prime factorized as $2^2 \cdot 5 \cdot 101$, it has $(2+1)(1+1)(1+1) = 12$ factors. So if we sum all the digits of all possible $a$ values, and add $4 \cdot 12 = 48$, we obtain the answer.

 Now we list out all factors of $2,020$, or all possible values of $a$. $1,2,4,5,10,20,101,202,404,505,1010,2020$. If we add up these digits, we get $45$, for a final answer of $45+48=\boxed{093}$.

 -molocyxu

 





\textbf{Solution 2 (Official MAA)}

Suppose that $N$ has the required property. Then there are positive integers $k$ and $m$ such that $N = 10^4m + 2020 = k\cdot m$. Thus $(k - 10^4)m = 2020$, which holds exactly when $m$ is a positive divisor of $2020.$ The number $2020 = 2^2\cdot 5\cdot 101$ has $12$ divisors: $1, 2, 4, 5, 10, 20, 101, 202, 404, 505, 1010$, and $2020.$ The requested sum is therefore the sum of the digits in these divisors plus $12$ times the sum of the digits in $2020,$ which is
 \[(1+2+4+5+1+2+2+4+8+10+2+4)+12\cdot4 = 93.\]

 





\textbf{Solution 3}

Note that for all $N \in S$, $N$ can be written as $N=10000x+2020=20(500x+101)$ for some positive integer $x$. Because $N$ must be divisible by $x$, $\frac{20(500x+101)}{x}$ is an integer. We now let $x=ab$, where $a$ is a divisor of $20$. Then $\frac{20(500x+101)}{x}=(\frac{20}{a})( \frac{500x}{b}+\frac{101}{b})$. We know $\frac{20}{a}$ and $\frac{500x}{b}$ are integers, so for $N$ to be an integer, $\frac{101}{b}$ must be an integer. For this to happen, $b$ must be a divisor of $101$. $101$ is prime, so $b\in \left \{ 1, 101 \right \}$. Because $a$ is a divisor of $20$, $a \in \left \{ 1,2,4,5,10,20\right\}$. So $x \in \left\{1,2,4,5,10,20,101,202,404,505,1010,2020\right\}$. Be know that all $N$ end in $2020$, so the sum of the digits of each $N$ is the sum of the digits of each $x$ plus $2+0+2+0=4$. Hence the sum of all of the digits of the numbers in $S$ is $12 \cdot 4 +45=\boxed{093}$.

 





\textbf{Solution 4 (Juxtaposing )}

The property of all numbers will fall under the category $\mathcal{N}2020$ where $\mathcal{N}$ is a sequence of digits. For the result to be generalized it follows that $\mathcal{N}$ yields the property, through \href{https://artofproblemsolving.com/wiki/index.php?title=Juxtaposition&action=edit&redlink=1}{juxtaposition}, that no matter the order we may always remove a sequence of 2, 0, 2, and 0 such that the resulting $\mathcal{N}$ divides the number as a whole, which is the same as saying all values of $\mathcal{N}$ must divide the term 2020. Thus, all values of $\mathcal{N}$ fall into the divisors of 2020. Continue as in previous solutions to obtain $\boxed{093}$.

 \textbf{Juxtaposition} is not an official article on AoPS, and so I shall explain it here. \textbf{Juxtaposition} is the property that regards two terms, essentially identifying that both terms may be swapped or interacted with that does not change the overall meaning of the term. The most simple cases of \textbf{Juxtaposition} include the \href{https://artofproblemsolving.com/wiki/index.php?title=Commutative\_property}{commutative property}, which states $a+b = b+a$, the \href{https://artofproblemsolving.com/wiki/index.php?title=Associative\_property}{associative property}, and the \href{https://artofproblemsolving.com/wiki/index.php?title=Identity\_property&action=edit&redlink=1}{identity property} of arithmetic. In this definition (not to be confused by convention, as there usually isn't a way to define a juxtaposed object as it is rarely used in math comps! I just so happened to find this one :D), define a juxtaposed orientation of an object $O$ to be $O^\mathcal{x}$.

 Here, \textbf{Juxtaposition} is directly used in terms of altering the digits, showing that we may form a specific number $O^\mathcal{x}$ and remove the digits 2, 0, 2, 0 in the same order, and the resulting number must also divide the original number $O$ given that $\mathcal{N}$ is not altered. This forms a bijection, and one can then show by special usage of juxtapositions that a special orientation forces $\mathcal{N}$ as a factor of 2020, hence proven. $\square$

 
                \begin{itemize}
                    \item \href{https://youtu.be/5b9Nw4bQt\_k}{https://youtu.be/5b9Nw4bQt\_k}\item \href{https://youtu.be/djWzRC-jGYw}{https://youtu.be/djWzRC-jGYw}
                \end{itemize}


                




	\item (\textit{AIME I 2020}) Six cards numbered $1$ through $6$ are to be lined up in a row. Find the number of arrangements of these six cards where one of the cards can be removed leaving the remaining five cards in either ascending or descending order.


 
 \textbf{Answer: }052



\textbf{Solution 1}

Realize that any sequence that works (ascending) can be reversed for descending, so we can just take the amount of sequences that satisfy the ascending condition and multiply by two.

 If we choose any of the numbers $1$ through $6$, there are five other spots to put them, so we get $6 \cdot 5 = 30$. However, we overcount some cases. Take the example of $132456$. We overcount this case because we can remove the $3$ or the $2$. Therefore, any cases with two adjacent numbers swapped is overcounted, so we subtract $5$ cases (namely, $213456, 132456, 124356, 123546, 123465$,) to get $30-5=25$, but we have to add back one more for the original case, $123456$. Therefore, there are $26$ cases. Multiplying by $2$ gives the desired answer, $\boxed{052}$.

 





\textbf{Solution 2 (Inspired by 2018 CMIMC Combo Round)}

Similar to above, a $1-1$ correspondence between ascending and descending is established by subtracting each number from $7$.

 We note that the given condition is equivalent to "cycling" $123456$ for a contiguous subset of it. For example, 

 $12(345)6 \rightarrow 125346, 124536$

 It's not hard to see that no overcount is possible, and that the cycle is either $1$ "right" or $1$ "left." Therefore, we consider how many elements we flip by. If we flip $1$ or $2$ such elements, then there is one way to cycle them. Otherwise, we have $2$ ways. Therefore, the total number of ascending is $1 + 5 + 2(4 + 3 + 2 + 1) = 26$, and multiplying by two gives $\boxed{052}.$ 

 





\textbf{Solution 3 (Table)}

We create a table with each digit in a corresponding slot. A sample one is shown below.

 
                    \begin{figure}[H]
                        \centering
                        \includegraphics[width=0.5\textwidth]{Solution/AIME_2020_I_5_Solution_Figure_1.png}
                    \end{figure}
                    The values highlighted in red are repeated values of an element that is highlighted green in the table. We count the values highlighted green. There are \( 26\) of these values.

 Notice (like said in Solution 1) that there is symmetry between the ascending and descending placements of the cards. This means we must multiply \( 26\) by \( 2\), giving us \( 26 \times 2 =\) $\boxed{052}.$

 





\textbf{Solution 4}

Similarly to Solution 2, we find the number of ascending arrangements and multiply by 2. 

 We can choose $5$ cards to be the ascending cards, therefore leaving $6$ places to place the remaining card. There are $\binom{6}{5}\cdot 6=36$ to do this. However, since the problem is asking for the number of arrangements, we overcount cases such as $123456$. Notice that the only arrangements that overcount are $123456$ (case 1) or if two adjacent numbers of $123456$ are switched (case 2).

 $\text{Case 1: }$ This arrangement is counted $6$ times. Each time it is counted for any of the $5$ numbers selected. Therefore we need to subtract $5$ cases of overcounting.

 $\text{Case 2: }$ Each time $2$ adjacent numbers of switched, there is one overcount. For example, if we have $213456$, both $1$ or $2$ could be removed. Since there are $5$ possible switches, we need to subtract $5$ cases of overcounting.

 Therefore, we have $36-5-5=26$ total arrangements of ascending numbers. We multiply by two (for descending) to get the answer of $\boxed{052}.$ 

 





\textbf{Solution 5 (No Overcounting)}

Like in previous solutions, we will count the number of ascending arrangements and multiply by 2. 

 First, consider the arrangement 1-2-3-4-5-6. That gives us 1 arrangement which works. 

 Next, we can switch two adjacent cards. There are 5 ways to pick two adjacent cards, so this gives us 5 arrangements.

 Now, we can "cycle" 3 adjacent cards. For example, 1-2-3 becomes 2-3-1 which becomes 3-1-2. There are 4 ways to pick a set of 3 adjacent cards, so this gives us 4x2=8 arrangements.

 Cycling 4 adjacent cards, we get the new arrangements 2-3-4-1 (which works), 3-4-1-2 (which doesn't work), and 4-1-2-3 (which does work). We get 6 arrangements.

 Similarly, when cycling 5 cards, we find 2x2=4 arrangements, and when cycling 6 cards, we find 2x1=2 arrangements.

 Adding, we figure out that there are 1+5+8+6+4+2=26 ascending arrangements. Multiplying by 2, we get the answer $\boxed{052}.$ 

 





\textbf{Solution 6 (Official MAA 1)}

First count the number of permutations of the cards such that if one card is removed, the remaining cards will be in ascending order. There is $1$ such permutation where all the cards appear in order: $123456.$ There are $5$ such permutations where two adjacent cards are interchanged, as in $124356.$ The other such permutations arise from removing one card from $123456$ and placing it in a position at least two away from its starting location. There are $4$ such positions to place each of the cards numbered $1$ and $6,$ and $3$ such positions for each of the cards numbered $2, 3, 4,$ and $5.$ This accounts for $2\cdot4 + 4\cdot3 =20$ permutations. Thus there are $1 + 5 + 20 = 26$ permutations where one card can be removed so that the remaining cards are in ascending order. There is an equal number of permutations that result in the cards' being in descending order. This gives the total $26 + 26 = \boxed{52}$.

 





\textbf{Solution 7 (Official MAA 2)}

More generally, suppose there are $n \geq 4$ cards numbered $1, 2, 3, \dots, n$ arranged in ascending order. If any one of the $n$ cards is removed and placed in one of the $n$ positions in the arrangement, the resulting permutation will have the property that one card can be removed so that the remaining cards are in ascending order. This accounts for $n\cdot n = n^2$ permutations. However, the original ascending order has been counted $n$ times, and each order that arises by switching two neighboring cards has been counted twice. Hence the number of arrangements where one card can be removed resulting in the remaining cards' being in ascending order is $n^2-(n-1)-(n-1)=(n-1)^2+1.$ When $n = 6$, this is $(6-1)^2+1 = 26$, and the final answer is $2\cdot26 = \boxed{52}$.

 





\textbf{Solution 8 (Casework)}

For ascending, if the $1$ goes in anything but the first two slots, the rest of the numbers have to go in ascending from $2$, which are $4$ cases if there are $6$ cards. If $1$ goes in the second spot, then you can put any of the rest in the first slot but then the rest are determined, so in the case of $6$ cards, that gives $5$ more. If $1$ goes in the first slot, that means that you are doing the same problem with $n-1$ cards. So the recursion is $a_n=(n-2)+(n-1)+a_{n-1}$. There's $a_1=1$ and $a_2=2$, so you get $a_3=2+3=5$, $a_4=5+5=10$, $a_5=7+10=17$, and $a_6=9+17=26$. Or you can see that $a_n=(n-1)^2+1$. We double to account for descending and get $\boxed{052}$. 

 





\textbf{Solution 9 (Symmetry and Case Study)}

First, we know that ascending order and descending order are symmetrical to each other (namely, if we get 132456 where after we take out 3, it will be one scenario; and if we flip it and write 654231, it will be another scenario)

 Thus, we only need to consider either descending or ascending and then times 2.WLOG let us consider ascending order

 Case 1: after we take out 1, the rest will be in ascending order:

 2 3 4 5 6Notice that 1 can be tucked in any one of the 6 spaces, thus there are 6 scenarios.

 Case 2: after we take out 2, the rest will be in ascending order:

 1 3 4 5 6Notice that if we put 2 next to 1 (to the right or to the left of 1), it will be an overcount, so there are only 4 cases for 2.

 It is easy to see that this is the same for 3, 4, 5, and 6.

 Thus, in total, we have 
 \[(6+4\times5)\times2=\boxed{052}\]~Adali

 





\textbf{Solution 10 (Illustration)}

We start with five cards in ascending order, then insert the sixth card to obtain a valid arrangement.

 Based on the card to be inserted, we have six cases. As shown below, the red squares indicate the possible positions to insert the sixth card.
                    \begin{figure}[H]
                        \centering
                        \includegraphics[width=0.5\textwidth]{Solution/AIME_2020_I_5_Solution_Figure_1.png}
                    \end{figure}
                    Note that all arrangements of the six cards are distinct.

 There are $26$ arrangements in which removing one card will leave the remaining five cards in ascending order. By symmetry, there are $26$ arrangements in which removing one card will leave the remaining five cards in descending order. So, the answer is $26\cdot2=\boxed{052}.$

 





\textbf{Solution 11}

We first realize that as long as we have an ascending sequence of $5$ numbers, we can just plug in a $6$th to make a sequence that works. For example, if we have $12345$, we can plug in a $6$ in any of $6$ spaces, before $1$, between $1$ and $2$, and so on to after $5$. We can also realize that this is completely symmetrical if the sequence is in descending order. For example, we could have $54321$, and we could plug in a 6 in 6 of the spaces. For the total number of combinations, we have 6 ascending cases multiplied by the $6$ places that can hold whatever number is missing, and multiply by 2 because there is descending and ascending number cases. 

 But wait, what if we have $12345$ and plug in a $6$ at the end and $12346$ and plug in a $5$ at the $5$th spot? We have overcounted, so we need to subtract off the identical pairs. Assume that one of them is the "right" combination. That means that there are $1$ "right" combination because all the rest will have $2$ of their combinations taken by the previous one. For example, if I have $12356$ and $12346$, I would count those $123456$ as $12356$'s instead of $12346$'s. Therefore, we have $6+4+4+4+4+4 = 26$ combinations of ascending order, and since we need to count those of descending order as well, we have $26\cdot2=\boxed{052}.$

 Side note: We could have just realized that the first part was unnecessary, but it does help you understand the problem better.

 





\textbf{Solution 12}

Like the previous solutions, we calculate ascending possibilities then multiply by two. Then, the configuration can look like +++++ or +-+++ (the minus can go anywhere), where "+" indicates an increase, and "-" indicates a decrease. Obviously, for the first case, the only possibility is 123456. For the second case, we can start with 123456, and then take either a card after the "-" and put it before the "-", or we could take a card before the "-" and put it after the "-". This means that every card other than the one at the position of the "-" can be used in one way. There are five ways to have 4 plus's, and one minus, and for each way, there are 5 ways to rearrange 123456 to achieve that. So, our answer is $2(5\cdot5+1)=\boxed{052}$

 





\textbf{Solution 13 (Undercount+add cases)}

First, we select $5$ of the $6$ cards to put in ascending or descending order. Then, we must add the other card to the group to arrange the six cards such that one can be removed so that the remaining cards are in ascending or descending order. Then, it would seem that the answer is $(6C5)*2$(ascending or descending)*$6$(the number of ways to place the last card). However, this overcounts some cases. For example, 123456 and 654321 can be made with 12345 and 6, 12346 and 5, etc. Furthermore, 654312 can be made with 65432 and 1 or 65431 and 2. Therefore, we first find that $(6C5)*2*4$=$48$. Then, we add the special cases ($123456$, $654321$, $654312$, $123465$) to obtain $\boxed{052}$.

 




	\item (\textit{AIME I 2020}) A flat board has a circular hole with radius $1$ and a circular hole with radius $2$ such that the distance between the centers of the two holes is $7$. Two spheres with equal radii sit in the two holes such that the spheres are tangent to each other. The square of the radius of the spheres is $\frac{m}{n}$, where $m$ and $n$ are relatively prime positive integers. Find $m+n$.


 
 \textbf{Answer: }173



\textbf{Solution 1 (Pythagorean Theorem)}


                    \begin{figure}[H]
                        \centering
                        \includegraphics[width=0.5\textwidth]{Solution/AIME_2020_I_6_Solution_Figure_1.png}
                    \end{figure}
                    Set the common radius to $r$. First, take the cross section of the sphere sitting in the hole of radius $1$. If we draw the perpendicular bisector of the chord (the hole) through the circle, this line goes through the center. Connect the center also to where the chord hits the circle, for a right triangle with hypotenuse $r$ and base $1$. Therefore, the height of this circle outside of the hole is $\sqrt{r^2-1}$.

 The other circle follows similarly for a height (outside the hole) of $\sqrt{r^2-4}$. Now, if we take the cross section of the entire board, essentially making it two-dimensional, we can connect the centers of the two spheres, then form another right triangle with base $7$, as given by the problem. The height of this triangle is the difference between the heights of the parts of the two spheres outside the holes, which is $\sqrt{r^2-1} - \sqrt{r^2-4}$. Now we can set up an equation in terms of $r$ with the Pythagorean theorem: 
 \[\left(\sqrt{r^2-1} - \sqrt{r^2-4}\right)^2 + 7^2 = (2r)^2.\] Simplifying a few times,
 \begin{align*} r^2 - 1 - 2\left(\sqrt{(r^2-1)(r^2-4)}\right) + r^2 - 4 + 49 &= 4r^2 \\ 2r^2-44 &= -2\left(\sqrt{(r^2-1)(r^2-4)}\right) \\ 22-r^2 &= \left(\sqrt{r^4 - 5r^2 + 4}\right) \\ r^4 -44r^2 + 484 &= r^4 - 5r^2 + 4 \\ 39r^2&=480 \\ r^2&=\frac{480}{39} = \frac{160}{13}. \end{align*}
Therefore, our answer is $\boxed{173}$.

 





\textbf{Solution 2  (Tangential Distance)}


                    \begin{figure}[H]
                        \centering
                        \includegraphics[width=0.5\textwidth]{Solution/AIME_2020_I_6_Solution_Figure_1.png}
                    \end{figure}
                    Let $A$ and $B$ be the centers of the holes, let $C$ be the point of crossing $AB$ and  radical axes of the circles. So $C$ has equal tangential distance to any point of both spheres. In particular to the circles (\href{https://en.wikipedia.org/wiki/Radical\_axis}{https://en.wikipedia.org/wiki/Radical\_axis}.)

 
                    \begin{figure}[H]
                        \centering
                        \includegraphics[width=0.5\textwidth]{Solution/AIME_2020_I_6_Solution_Figure_1.png}
                    \end{figure}
                    
 \[CA = \frac {AB} {2} – \frac {r_A^2 – r_B^2}{2 AB} = \frac{23}{7}, CB = AB - AC =\frac{26}{7},\]
 \[CA' = CB'= \sqrt{BC^2 – r_B^2} = \frac {4}{7} \sqrt{30}.\]

 Let $D$ be the point of tangency of the spheres common radius $R$ centered at $O$ and $O'.$ Let $\alpha$ be the angle between $OO'$ and flat board. In the plane, perpendicular to board 
 \[DC \perp OO', DC =  \frac {4}{7} \sqrt{30}.\]

 
                    \begin{figure}[H]
                        \centering
                        \includegraphics[width=0.5\textwidth]{Solution/AIME_2020_I_6_Solution_Figure_1.png}
                    \end{figure}
                    Distance between $C$ and midpoint $M$ of $AB$ is 
 \[d = \frac {26}{7} - \frac {7}{2} = \frac{3}{14} \implies \sin \alpha = \frac {d}{DC} = \frac {3}{8\sqrt {30}}.\]
 \[\cos \alpha = \sqrt {1 – \frac {9}{64 \cdot 30}} = \sqrt{ \frac {637}{640}} = \frac {7}{2} \sqrt {\frac{13}{160}}.\]
 \[2R \cos \alpha = AB = 7 \implies R = \frac {\frac{7}{2} } {\frac{7}{2}\sqrt \frac{13}{160}} = \sqrt {\frac{160}{13}}.\]Hence the answer is 
 \[160+13 = \boxed{173}.\]\textbf{vladimir.shelomovskii@gmail.com, vvsss}

 




	\item (\textit{AIME I 2020}) A club consisting of $11$ men and $12$ women needs to choose a committee from among its members so that the number of women on the committee is one more than the number of men on the committee. The committee could have as few as $1$ member or as many as $23$ members. Let $N$ be the number of such committees that can be formed. Find the sum of the prime numbers that divide $N.$


 
 \textbf{Answer: }081



\textbf{Solution 1}

Let $k$ be the number of women selected. Then, the number of men not selected is $11-(k-1)=12-k$.Note that the sum of the number of women selected and the number of men not selected is constant at $12$. Each combination of women selected and men not selected corresponds to a committee selection. Since choosing 12 individuals from the total of 23 would give $k$ women and $12-k$ men, the number of committee selections is $\binom{23}{12}$. The answer is $\boxed{081}$.~awang11's sol

 





\textbf{Solution 2 (Bash)}

We casework on the amount of men on the committee. 

 If there are no men in the committee, there are $\dbinom{12}{1}$ ways to pick the women on the committee, for a total of $\dbinom{11}{0} \cdot \dbinom{12}{1}$. Notice that $\dbinom{11}{0}$ is equal to $\dbinom{11}{11}$, so the case where no men are picked can be grouped with the case where all men are picked. When all men are picked, all women must also be picked, for a total of $\dbinom{12}{12}$. Therefore, these cases can be combined to 
 \[\dbinom{11}{0} \cdot \left(\dbinom{12}{1} + \dbinom{12}{12}\right)\] Since $\dbinom{12}{12} = \dbinom{12}{0}$, and $\dbinom{12}{0} + \dbinom{12}{1} = \dbinom{13}{1}$, we can further simplify this to 
 \[\dbinom{11}{0} \cdot \dbinom{13}{1}\]

 All other cases proceed similarly. For example, the case with one men or ten men is equal to $\dbinom{11}{1} \cdot \dbinom{13}{2}$. Now, if we factor out a $13$, then all cases except the first two have a factor of $121$, so we can factor this out too to make our computation slightly easier. The first two cases (with $13$ factored out) give $1+66=67$, and the rest gives $121(10+75+270+504) = 103,939$. Adding the $67$ gives $104,006$. Now, we can test for prime factors. We know there is a factor of $2$, and the rest is $52,003$. We can also factor out a $7$, for $7,429$, and the rest is $17 \cdot 19 \cdot 23$. Adding up all the prime factors gives $2+7+13+17+19+23 = \boxed{081}$.

 \textbf{Video Solution:} 





\textbf{Solution 3 (Vandermonde's identity)}

Applying \href{https://artofproblemsolving.com/wiki/index.php/Combinatorial\_identity}{[1]} by setting $m=12$, $n=11$, and $r=11$, we obtain $\binom{23}{11}\implies$ $\boxed{081}$.~Lcz

 \textbf{Short Proof} Consider the following setup:
                    \begin{figure}[H]
                        \centering
                        \includegraphics[width=0.5\textwidth]{Solution/AIME_2020_I_7_Solution_Figure_1.png}
                    \end{figure}
                    The dots to the left represent the men, and the dots to the right represent the women. Now, suppose we put a mark on $11$ people (the $*$). Those to the left of the dashed line get to be "in" on the committee if they have a mark. Those on the right side of the dashed line are already on the committee, but if they're marked they get forcibly evicted from it. If there were $x$ people marked on the left, there ends up being $12-(11-x) = x+1$ people not marked on the right. Circles represent those in the committee.
                    \begin{figure}[H]
                        \centering
                        \includegraphics[width=0.5\textwidth]{Solution/AIME_2020_I_7_Solution_Figure_1.png}
                    \end{figure}
                    

 We have our bijection, so the number of ways will be $\binom{23}{11}$.

 





\textbf{Solution 4}

Notice that the committee can consist of $k$ boys and $k+1$ girls. Summing over all possible $k$ gives 
 \[\sum_{k=0}^{11}\binom{11}{k}\binom{12}{k+1}=\binom{11}{0}\binom{12}{1}+\binom{11}{1}\binom{12}{2}+\cdots + \binom{11}{11}\binom{12}{12}\]Using the identity $\binom{n}{k}=\binom{n}{n-k}$, and Pascal's Identity $\binom{n}{k}+\binom{n}{k+1}=\binom{n+1}{k+1}$, we get
 \[\sum_{k=0}^{11}\binom{11}{k}\binom{12}{k+1}=\binom{12}{12}+\binom{12}{1}\left(\binom{11}{0}+\binom{11}{1}\right)+\cdots\]
 \[=\binom{12}{0}^2+\binom{12}{1}^2+\binom{12}{2}^2+\binom{12}{3}^2+\binom{12}{4}^2+\binom{12}{5}^2+\frac{\binom{12}{6}^2}{2}\]
 \[=\frac{1}{2}\sum_{k=0}^{12}\binom{12}{k}^2\]Using the identity $\sum_{k=0}^n\binom{n}{k}^2=\binom{2n}{n}$, this simplifies to 
 \[\frac{1}{2}\cdot \binom{24}{12}=\frac{24\cdot 23\cdot 22\cdot 21\cdot 20\cdot 19\cdot 18\cdot 17\cdot 16\cdot 15\cdot 14\cdot 13}{2\cdot 12\cdot 11\cdot 10\cdot 9\cdot 8\cdot 7\cdot 6\cdot 5\cdot 4\cdot 3\cdot 2}=2\cdot 7\cdot 13\cdot 17\cdot 19\cdot 23\]so the desired answer is $2+7+13+17+19+23=\boxed{081}$ ~ktong

 





\textbf{Solution 5 (Official MAA)}

Select any $11$ club members. That group will have $i$ men and $11-i$ women, so the number of women in the club not selected in that group is $12 - (11-i) = i+1$. Thus, if the committee includes the men who were selected and the women who were not selected, the committee would have the correct number of men and women. Conversely, for every committee that could be formed with $i$ men and $i+1$ women, the men on this committee together with the women not on the committee comprise a subset of $i + (12 - (i+1)) = 11$ club members. Thus
 \[N = \binom{23}{11}= \frac{23\cdot22\cdot21\cdot20\cdot19\cdot18\cdot17\cdot16\cdot15\cdot14\cdot13}{11\cdot10\cdot9\cdot8\cdot7\cdot6\cdot5\cdot4\cdot3\cdot2\cdot1}=23\cdot19\cdot17\cdot13\cdot7\cdot2.\]The requested sum is $23+19+17+13+7+2=\boxed{081}$.

 





\textbf{Solution 6}

Notice that if $k$ men are picked, then $k+1$ women must be picked. Furthermore, $k$ can range from $0$ to $11$. Then, 
 \[N=\sum_{k=0}^{11} \binom{11}{k}\binom{12}{k+1}=\binom{11}{0}\binom{12}{1}+\binom{11}{1}\binom{12}{2}+\dots+\binom{11}{11}\binom{12}{12}\]Since $\binom{n}{k}=\binom{n}{n-k}$, this equals
 \[N=\binom{11}{0}\binom{12}{11}+\binom{11}{1}\binom{12}{10}+\dots+\binom{11}{11}\binom{12}{0}\]According to Vandermonde's Identity,
 \[N=\binom{11+12}{11}=\binom{23}{11}\]
 \[N=\frac{23!}{11!\cdot 12!}=\frac{2^{19}\cdot 3^9\cdot 5^4\cdot 7^3\cdot 11^2\cdot 13\cdot 17\cdot 19\cdot 23}{2^{10}\cdot 3^5\cdot 5^2\cdot 7\cdot 11\times 2^8\cdot 3^4\cdot 5^2\cdot 7\cdot 11}=2\cdot 7\cdot 13\cdot 17\cdot 19\cdot 23\rightarrow \boxed{081}\]~\href{https://artofproblemsolving.com/wiki/index.php/User:Sid2012\#User}{sid2012}

 





\textbf{Solution 7 ( *Rigorous* Recursion)}

Test the cases where there are $0$ men $1$ woman, $1$ man $2$ women, $2$ men $3$ women ... you will get the sequence $1$, $3$, $10$, $35$. Multiply all these numbers by $2$ to get $2$, $6$, $20$, $70$, which is also $\binom{2}{1}$, $\binom{4}{2}$, $\binom{6}{3}$, $\binom{8}{4}$. Thus, continuing this pattern, the case with $11$ men and $12$ women should have $\frac{\binom{24}{12}}{2}$ ways to select the committee.-Kevin2010

 (Solves using both methods - Casework and Vandermonde's Identity)

 




	\item (\textit{AIME I 2020}) A bug walks all day and sleeps all night. On the first day, it starts at point $O,$ faces east, and walks a distance of $5$ units due east. Each night the bug rotates $60^\circ$ counterclockwise. Each day it walks in this new direction half as far as it walked the previous day. The bug gets arbitrarily close to the point $P.$ Then $OP^2=\frac{m}{n},$ where $m$ and $n$ are relatively prime positive integers. Find $m+n.$


 
 \textbf{Answer: }103



\textbf{Solution 1}

Put the problem on the complex plane. Let $O$ be $0$. Define $\omega := \frac{1}{2}e^{\frac{\pi i}{3}} =\frac{1}{4}+\frac{\sqrt{3}}{4}i$. The bug's original movement vector is the complex number $5$, and each day this vector is multiplied by $\omega$, so summing up all these vectors, we find
 \[OP^2 = \left| 5+5\omega + 5\omega^2 + \dots \right|^2 = \left| \frac{5}{1-\omega} \right|^2 = \frac{25}{|1-\omega|^2} = \frac{25}{|\frac{3}{4}-\frac{\sqrt{3}}{4}i|^2}=\frac{25}{\frac{3}{4}}=\frac{100}{3} \implies \boxed{103}\]

 -sillybone

 





\textbf{Solution 2 (Coordinates)}


                    \begin{figure}[H]
                        \centering
                        \includegraphics[width=0.5\textwidth]{Solution/AIME_2020_I_8_Solution_Figure_1.png}
                    \end{figure}
                    

 We notice that the moves cycle every 6 moves, so we plot out the first 6 moves on the coordinate grid with point $O$ as the origin. We will keep a tally of the x-coordinate and y-coordinate separately. Then, we will combine them and account for the cycling using the formula for an infinite geometric series.

 First move: The ant moves right $5$. Second move: We use properties of a $30-60-90$ triangle to get $\frac{5}{4}$ right, $\frac{5\sqrt{3}}{4}$ up. Third move: $\frac{5}{8}$ left, $\frac{5\sqrt{3}}{8}$ up.Fourth move: $\frac{5}{8}$ left.Fifth move: $\frac{5}{32}$ left, $\frac{5\sqrt{3}}{32}$ down.Sixth move: $\frac{5}{64}$ right, $\frac{5\sqrt{3}}{64}$ down.

 Total of x-coordinate: $5 + \frac{5}{4} - \frac{5}{8} - \frac{5}{8}  - \frac{5}{32} + \frac{5}{64} = \frac{315}{64}$.Total of y-coordinate: $0 + \frac{5\sqrt{3}}{4} + \frac{5\sqrt{3}}{8} + 0 - \frac{5\sqrt{3}}{32} - \frac{5\sqrt{3}}{64} = \frac{105\sqrt{3}}{64}$.

 After this cycle of six moves, all moves repeat with a factor of $(\frac{1}{2})^6 = \frac{1}{64}$. Using the formula for a geometric series, multiplying each sequence by $\frac{1}{1-\frac{1}{64}} = \frac{64}{63}$ will give us the point $P$.

 Now, knowing the initial $x$ and $y,$ we plug this into the geometric series formula ($\frac{a}{1-r}$), and we get $\frac{315}{64} \cdot \frac{64}{63} = 5$, $\frac{105\sqrt{3}}{64} \cdot \frac{64}{63} = \frac{5\sqrt{3}}{3}$.Therefore, the coordinates of point $P$ are $(5,\frac{5\sqrt{3}}{3})$, so using the Pythagorean Theorem, $OP^2 = \frac{100}{3}$, for an answer of $\boxed{103}$.

 
-molocyxu

 





\textbf{Solution 3 (Complex)}

We place the ant at the origin of the complex plane with its first move being in the positive real direction.  Then the ant's journey can be represented as the infinite series
 \[5\left(1 + \frac{e^{\frac{i\pi}{3}}}{2} + \left(\frac{e^{\frac{i\pi}{3}}}{2}\right)^2 + \cdots\right)\]Using the formula for an infinite geometric series, this is equal to
 \[\frac{5}{1 - \frac12e^{\frac{i\pi}{3}}} = \frac{5}{1 - \frac{1 + i\sqrt{3}}{4}} = \frac{20}{3 - i\sqrt{3}} = 5 + \frac{5i\sqrt{3}}{3}\]We are looking for the square of the modulus of this value:
 \[\left|5 + \frac{5i\sqrt{3}}{3}\right|^2 = 25 + \frac{25}{3} = \frac{100}{3}\]so the answer is $100 + 3 = \boxed{103}$.

 





\textbf{Solution 3 (Solution 1 Faster)}

The ant goes in the opposite direction every $3$ moves, going $(1/2)^3=1/8$ the distance backwards. Using geometric series, he travels $1-1/8+1/64-1/512...=(7/8)(1+1/64+1/4096...)=(7/8)(64/63)=8/9$ the distance of the first three moves over infinity moves. Now, we use coordinates meaning $(5+5/4-5/8, 0+5\sqrt3/4+5\sqrt3/8)$ or $(45/8, 15\sqrt3/8)$. Multiplying these by $8/9$, we get $(5, 5\sqrt3/3)$ $\implies$ $\boxed{103}$ .

 





\textbf{Solution 5 (Official MAA 1)}

Suppose that the bug starts at the origin $(0,0)$ and travels a distance of $a$ units due east on the first day, and that there is a real number $r$ with $0<r < 1$ such that each day after the first, the bug walks $r$ times as far as the previous day. On day $n$, the bug travels along the vector $\pmb v_{n}$ that has magnitude $ar^{n-1}$ and direction $\langle\cos(n\cdot 60^\circ),\sin(n\cdot 60^\circ)\rangle$. Then $P$ is the terminal point of the infinite sum of the vectors $\pmb v_{1}+\pmb v_{2}+\pmb v_3+\cdots$. The $x$-coordinate of this sum is
 \[a\big(\!\cos0^\circ+r\cos60^\circ + r^{2}\cos120^\circ+r^{3}\cos180^\circ+r^{4}\cos240^\circ\]
 \[+r^{5}\cos300^\circ+r^{6}\cos360^\circ+\cdots\big).\]Because the angles repeat after 6 terms, this sum is equal to
 \[aS(1+r^{6}+r^{12}+r^{18}+\cdots)=\frac{aS}{1-r^{6}},\]where
 \[S=\cos0^\circ+r\cos60^\circ + r^{2}\cos120^\circ+r^{3}\cos180^\circ+r^{4}\cos240^\circ+ r^{5}\cos300^\circ.\]Similarly, the $y$-coordinate of $P$ will be ${\frac{aT}{1-r^{6}}}$, where
 \[T=\sin0^\circ+r\sin60^\circ + r^{2}\sin120^\circ+r^{3}\sin180^\circ+r^{4}\sin240^\circ+ r^{5}\sin300^\circ.\]In this case $r=\frac12$ and $a = 5$, so
 \[S=1+\frac14-\frac18-\frac18-\frac1{32}+\frac1{64}=\frac{63}{64},\]
 \[T=0+\frac{\sqrt3}4+\frac{\sqrt3}8+0-\frac{\sqrt3}{32}-\frac{\sqrt3}{64}=\frac{21\sqrt3}{64},\]and the coordinates of $P$ are
 \[\left(\frac{5S}{1-\frac1{64}}, \frac{5T}{1-\frac1{64}}\right)=\left(5,\frac{5\sqrt3}{3}\right).\]Thus the square of the distance from the origin to $P$ is $25+\frac{25}3=\frac{100}3$. The requested sum is $100+3=\boxed{103}$.

 





\textbf{Solution 6 (Official MAA 2)}

Let point $O$ be the origin in the complex plane. Point $P$ is the complex sum $5(1+z+z^2+\cdots) = \frac{5}{1-z}$, where $z=\frac{1+i\sqrt3}4$. The distance squared is
 \[{OP}^2=\left|\frac5{1-\frac{1+i\sqrt3}4}\right|^{2}= \frac{(4\cdot5)^2}{\left|4-(1+i\sqrt3)\right|^2}=\frac{400}{9+3}=\frac{100}3.\]Hence the answer is $100 + 3 = \boxed{103}$.

 





\textbf{Solution 7 (No coordinates, but basically using the same idea as Solution 1)}

The bug goes forward and backward in three directions: straight east, northeast, and northwest. It travels $5-\frac{5}{8}+\frac{5}{64}-\cdots=\frac{40}{9}$ units east. Thus, it goes northeast $\frac{\frac{40}{9}}{2}=\frac{20}{9}$ units northeast and $\frac{10}{9}$ units northwest. Now, the bug travels a total of $\frac{40}{9}+\frac{10}{9}-\frac{5}{9}=\frac{45}{9}=5$ units east and a total of $\frac{10\sqrt{3}}{9}+\frac{5\sqrt{3}}{9}=\frac{15\sqrt{3}}{9}=\frac{5\sqrt{3}}{3}$ units north because of the 30-60-90 right triangles formed. Now, $OP^2=5^2+\frac{5^2}{3}=\frac{100}{3}$ by the Pythagorean Theorem, and the answer is $\boxed{103}$.

 -integralarefun

 





\textbf{Solution 8}

The bug's bearings on each traversal are $0^\circ, 60^\circ, 120^\circ,$ and so on; in general, the $n-$th traversal has length $5\cdot (1/2)^{n-1}$ and bearing $60(n-1).$ This means that the $x$ and $y$ displacements for the $n-$th traversal are

 
 \[(\Delta x_n, \Delta y_n)=(5\cdot (1/2)^{n-1}\cos (60(n-1))^\circ,5\cdot (1/2)^{n-1}\sin (60(n-1))^\circ).\]

 Summing this over all the displacements, we get

 
 \[x_P=\sum_{n=1}^{\infty} 5\cdot (1/2)^{n-1}\cos (60(n-1))^\circ, y_P=\sum_{n=1}^{\infty} 5\cdot (1/2)^{n-1}\sin (60(n-1))^\circ.\]

 We then have 

 
 \begin{align} OP^2 &= x_P^2+y_P^2 \\ &=\sum_{n=1}^{\infty} (5\cdot (1/2)^{n-1}\cos ^2(60(n-1))^\circ) + (5\cdot (1/2)^{n-1}\sin ^2(60(n-1))^\circ) \\ &= \sum_{n=1}^{\infty} (5\cdot (1/2)^{n-1})^2(\cos ^2 (60(n-1))^\circ+\sin ^2 (60(n-1))^\circ) \\ &= \sum_{n=1}^{\infty} (25\cdot (1/4)^{n-1}) \\ &= \dfrac{25}{1-1/4} \\ &= 100/3. \end{align}

 Thus, the answer is $100+3=\boxed{103}.$ --MenuThreeOne

 Moving Bug in AIME solution with python turtle scripts

 -avn

 




	\item (\textit{AIME I 2020}) Let $S$ be the set of positive integer divisors of $20^9.$ Three numbers are chosen independently and at random with replacement from the set $S$ and labeled $a_1,a_2,$ and $a_3$ in the order they are chosen. The probability that both $a_1$ divides $a_2$ and $a_2$ divides $a_3$ is $\frac{m}{n},$ where $m$ and $n$ are relatively prime positive integers. Find $m.$


 
 \textbf{Answer: }077



\textbf{Solution 1}


                    \begin{figure}[H]
                        \centering
                        \includegraphics[width=0.5\textwidth]{Solution/AIME_2020_I_9_Solution_Figure_1.png}
                    \end{figure}
                    

 First, prime factorize $20^9$ as $2^{18} \cdot 5^9$. Denote $a_1$ as $2^{b_1} \cdot 5^{c_1}$, $a_2$ as $2^{b_2} \cdot 5^{c_2}$, and $a_3$ as 

 




	\item (\textit{AIME I 2020}) Let $m$ and $n$ be positive integers satisfying the conditions


 $\quad\bullet\ \gcd(m+n,210)=1,$


 $\quad\bullet\ m^m$ is a multiple of $n^n,$ and


 $\quad\bullet\ m$ is not a multiple of $n.$


 Find the least possible value of $m+n.$


 
 \textbf{Answer: }407



\textbf{Solution 1}

Taking inspiration from $4^4 \mid 10^{10}$ we are inspired to take $n$ to be $p^2$, the lowest prime not dividing $210$, or $11 \implies n = 121$. Now, there are $242$ factors of $11$, so $11^{242} \mid m^m$, and then $m = 11k$ for $k \geq 22$. Now, $\gcd(m+n, 210) = \gcd(11+k,210) = 1$. Noting $k = 26$ is the minimal that satisfies this, we get $(n,m) = (121,286)$. Thus, it is easy to verify this is minimal and we get $\boxed{407}$. ~awang11

 





\textbf{Solution 2}

We essentially have that $\gcd(m,n) \notin \{1, m, n\},$ or the set of factors of $210.$

 
This implies that one number must be odd and the other must be even so that they don’t add up to a multiple of $2.$ If $n$ is even and $m$ is odd it would be impossible for $m^m$ to be a multiple of $n^n,$ so $m$ must be even and $n$ must be odd. 

 Now we need to choose the prime factor for $m$ and $n$ to share. It can't be $2,3,5,$ or $7$ so the next smallest option is $11.$ 

 So far we have $m = 22$ and $n = 11.$ We need to add more factors so that $m$ is not a multiple of $n$ but $m^m$ is a multiple of $n^n.$ Note that no matter how many additional factors we add to $m,$ it will always be a multiple of $11,$ so we have to add another factor to $n$ to make sure it doesn't divide $m.$ Again the smallest option is $11,$ so $n$ becomes $121 \implies n^n = 11^{242},$ and $m^m = 2^{22} \cdot 11^{22}.$ All we need to do now is significantly increase the value of $m$ so that the exponent on $11$ becomes larger than $242.$ If we added another factor of $11$ to $m$ then it would be a multiple of $n$ again, so the next smallest option is $13.$ Then $m = 22\cdot 13 = 286.$ $m^m$ becomes $2^{286} \cdot 11^{286} \cdot 13^{286}$ which satisfies the problem's condition.

 Therefore, the least possible value of $m + n$ is $121 + 286 = \boxed{407}.$ 

 





\textbf{Solution 3}

Assume for the sake of contradiction that $n$ is a multiple of a single digit prime number, then $m$ must also be a multiple of that single digit prime number to accommodate for $n^n | m^m$. However that means that $m+n$ is divisible by that single digit prime number, which violates $\gcd(m+n,210) = 1$, so contradiction.

 $n$ is also not 1 because then $m$ would be a multiple of it.

 Thus, $n$ is a multiple of 11 and/or 13 and/or 17 and/or...

 Assume for the sake of contradiction that $n$ has at most 1 power of 11, at most 1 power of 13...and so on...Then, for $n^n | m^m$ to be satisfied, $m$ must contain at least the same prime factors that $n$ has. This tells us that for the primes where $n$ has one power of, $m$ also has at least one power, and since this holds true for all the primes of $n$, $n|m$. Contradiction.

 Thus $n$ needs more than one power of some prime.The obvious smallest possible value of $n$ now is $11^2 =121$.Since $121^{121}=11^{242}$, we need $m$ to be a multiple of 11 at least $242$ that is not divisible by $121$ and most importantly, $\gcd(m+n,210) = 1$.$242$ is divisible by $121$, out.$253+121$ is divisible by 2, out.$264+121$ is divisible by 5, out.$275+121$ is divisible by 2, out.$286+121=37\cdot 11$ and satisfies all the conditions in the given problem, and the next case $n=169$ will give us at least $169\cdot 3$, so we get $\boxed{407}$.

 





\textbf{Solution 4 (Official MAA)}

Let $m$ and $n$ be positive integers where $m^m$ is a multiple of $n^n$ and $m$ is not a multiple of $n$. If a prime $p$ divides $n$, then $p$ divides $n^n$, so it also divides $m^m$, and thus $p$ divides $m$. Therefore any prime $p$ dividing $n$ also divides both $m$ and $k = m + n$. Because $k$ is relatively prime to $210=2\cdot3\cdot5\cdot7$, the primes $2$, $3$, $5$, and $7$ cannot divide $n$. Furthermore, because $m$ is divisible by every prime factor of $n$, but $m$ is not a multiple of $n$, the integer $n$ must be divisible by the square of some prime, and that prime must be at least $11$. Thus $n$ must be at least $11^2 = 121$.

 If $n=11^2$, then $m$ must be a multiple of $11$ but not a multiple of $121$, and $m^m$ must be divisible by $n^n = 121^{121} = 11^{242}$. Therefore $m$ must be a multiple of $11$ that is greater than $242$. Let $m = 11m_0$, with $m_0 > 22$. Then $k = m + n = 11(m_0 + 11)$. The least $m_0 > 22$ for which $m_0 + 11$ is not divisible by any of the primes $2$, $3$, $5$, or $7$ is $m_0 = 26$, giving the prime $m_0 + 11 = 37$. Hence the least possible $k$ when $n = 121$ is $k = 11 \cdot 37 = 407$.

 It remains to consider other possible values for $n$. If $n = 13^2 = 169$, then $m$ must be divisible by $13$ but not $169$, and $m^m$ must be a multiple of $n^n = 169^{169} = 13^{338}$, so $m > 338$. Then $k = m + n > 169 + 338 = 507$. All other possible values for $n$ have $n \ge 242$, and in this case $m > n \ge 242$, so $k \ge 2 \cdot 242 = 484$. Hence no greater values of $n$ can produce lesser values for $k$, and the least possible $k$ is indeed $407$.

 





\textbf{Solution 5}

Because $210 = 2\times3\times5\times7$, the minimum factor of $m+n$ can have is $11$.Suppose $(m,n) = 11$

 $\left\{ 	\begin{aligned} 		&m=11a\\ 		&n=11b 	\end{aligned} \right.$$(a,b) = 1$

 $11b \mid 11a$

 $\textbf{Min} b = 11$

 $n=11^2$

 So $a>22$

 We only need to test several more values of a to see whether the condition $\gcd(m+n,210)=1$ fits:

 $a=23, 2\mid a+b = 34$

 $a=24, 5\mid a+b = 35$

 $a=25, 2\mid a+b = 36$

 $a=26, a+b = 37$

 Therefore $\textbf{Min} (m+n) = 11(a+b) = 11\times37 = \boxed{407}$

 -cassphe

 


 




	\item (\textit{AIME I 2020}) For integers $a,b,c$ and $d,$ let $f(x)=x^2+ax+b$ and $g(x)=x^2+cx+d.$ Find the number of ordered triples $(a,b,c)$ of integers with absolute values not exceeding $10$ for which there is an integer $d$ such that $g(f(2))=g(f(4))=0.$


 
 \textbf{Answer: }510



\textbf{Solution 1 (Strategic Casework)}

There can be two different cases for this problem, either $f(2)=f(4)$ or not. If it is, note that by Vieta's formulas $a = -6$. Then, $b$ can be anything. However, $c$ can also be anything, as we can set the root of $g$ (not equal to $f(2) = f(4)$) to any integer, producing a possible integer value of $d$. Therefore there are $21^2 = 441$ in this case*. If it isn't, then $f(2),f(4)$ are the roots of $g$. This means by Vieta's, that:

 
 \[f(2)+f(4) = -c \in [-10,10]\]
 \[20 + 6a + 2b \in [-10,10]\]
 \[3a + b \in [-15,-5].\]

 Solving these inequalities while considering that $a \neq -6$ to prevent $f(2) = f(4)$, we obtain $69$ possible tuples and adding gives $441+69=\boxed{510}$. ~awang11

 





\textbf{Solution 2 (Bash)}

Define $h(x)=x^2+cx$. Since $g(f(2))=g(f(4))=0$, we know $h(f(2))=h(f(4))=-d$. Plugging in $f(x)$ into $h(x)$, we get $h(f(x))=x^4+2ax^3+(2b+a^2+c)x^2+(2ab+ac)x+(b^2+bc)$. Setting $h(f(2))=h(f(4))$, 
 \[16+16a+8b+4a^2+4ab+b^2+4c+2ac+bc=256+128a+32b+16a^2+8ab+b^2+16c+4ac+bc\]. Simplifying and cancelling terms, 
 \[240+112a+24b+12a^2+4ab+12c+2ac=0\] 
 \[120+56a+12b+6a^2+2ab+6c+ac=0\] 
 \[6a^2+2ab+ac+56a+12b+6c+120=0\] 
 \[6a^2+2ab+ac+20a+36a+12b+6c+120=0\] 
 \[a(6a+2b+c+20)+6(6a+2b+c+20)=0\] 
 \[(a+6)(6a+2b+c+20)=0\]

 Therefore, either $a+6=0$ or $6a+2b+c=-20$. The first case is easy: $a=-6$ and there are $441$ tuples in that case. In the second case, we simply perform casework on even values of $c$, to get $77$ tuples, subtracting the $8$ tuples in both cases we get $441+77-8=\boxed{510}$.

 -EZmath2006

 





\textbf{Solution 3 (Official MAA)}

For a given ordered triple $(a, b, c)$, the value of $g(f(4)) - g(f(2))$ is uniquely determined, and a value of $d$ can be found to give $g(f(2))$ any prescribed integer value. Hence the required condition can be satisfied provided that $a$, $b$, and $c$ are chosen so that
 \begin{align*} 0 &= g(f(4)) - g(f(2)) = \big(f(4)\big)^{\!2} - \big(f(2)\big)^{\!2} + c\big(f(4) - f(2)\big)\\ &= \big(f(4) - f(2)\big)\big(f(4) + f(2) + c\big)\\ &= (12 + 2a)(20 + 6a + 2b + c). \end{align*}
First suppose that $12 + 2a = 0$, so $a = -6$. In this case there are $21$ choices for each of $b$ and $c$ with $-10 \le b \le 10$ and $-10 \le c \le 10$, so this case accounts for $21^2 = 441$ ordered triples.

 Next suppose that $a \ne -6$ and $20 + 6a + 2b + c = 0$, so $c = -6a - 2b - 20$. Because$-10 \le c \le 10$, it follows that $-30 \le 6a + 2b \le -10$, and because $-10\le b \le10$, it follows that $-8 \le a \le 1$. Then$-15 - 3a \le b \le -5 - 3a$. The number of ordered triples for various values of $a$ are presented in the following table.

 
 \[\begin{tabular}{|c|c|c|r|} \hline a  & b                  & c & triples \\ \hline -8 & \{9,10\}           & -6a - 2b - 20    &             2              \\ -7 & \{6, 7, 8, 9, 10\} &   -6a - 2b - 20  &         5                  \\ -6 & \{-10, \ldots, 10\} &  \{-10, \ldots, 10\}   &                  441        \\ -5 & \{0, \ldots, 10\}      &  -6a - 2b - 20   &    11                       \\ -4 & \{-3, \ldots, 7\}        &  -6a - 2b - 20   &   11                       \\ -3 & \{-6, \ldots, 4\}        &  -6a - 2b - 20   &   11                       \\ -2 & \{-9, \ldots, 1\}         &  -6a - 2b - 20   &   11                        \\ -1 & \{-10, \ldots, -2\}   &  -6a - 2b - 20   &           9                \\ 0  & \{-10, \ldots, -5\}        &   -6a - 2b - 20  &      6                     \\ 1  & \{-10, -9, -8\}        &  -6a - 2b - 20   &      3                     \\ \hline Total & & & 510\\ \hline \end{tabular}\]The total number of ordered triples that satisfy the required condition is $\boxed{510}$.

 \textbf{Notes For *} In case anyone is confused by this (as I initially was). In the case where $f(2)=f(4)$, this does not mean that g has a double root of $f(2)=f(4)=k$, ONLY that $k$ is one of the roots of g. So basically since $a=-6$ in this case, $f(2)=f(4)=b-8$, and we have $21$ choices for b and we \textit{still can} ensure c is an integer with absolute value less than or equal to 10 simply by having another integer root of g that when added to $b-8$ ensures this, and of course an integer multiplied by an integer is an integer so $d$ will still be an integer. In other words, you have can have $b$ and $c$ be any integer with absolute value less than or equal to 10 with $d$ still being an integer. Now refer back to the 1st solution.~First

 




	\item (\textit{AIME I 2020}) Let $n$ be the least positive integer for which $149^n-2^n$ is divisible by $3^3\cdot5^5\cdot7^7.$ Find the number of positive integer divisors of $n.$


 
 \textbf{Answer: }270



\textbf{Solution 1.1}

As usual, denote $v_p(n)$ the highest power of prime $p$ that divides $n$.\href{https://artofproblemsolving.com/wiki/index.php?title=Lifting\_the\_Exponent}{Lifting the Exponent} shows that 
 \[3=v_3(149^n-2^n) = v_3(n) + v_3(147) = v_3(n)+1\] so thus, $3^2$ divides $n$. It also shows that 
 \[7=v_7(149^n-2^n) = v_7(n) + v_7(147) = v_7(n)+2\] so thus, $7^5$ divides $n$. 

 
Now, setting $n = 4c$ (necessitated by $149^n \equiv 2^n \pmod 5$ in order to set up LTE), we see 
 \[v_5(149^{4c}-2^{4c}) = v_5(149^{4c}-16^{c})\] and since $149^{4} \equiv 1 \pmod{25}$ and $16^1 \equiv 16 \pmod{25}$ then $v_5(149^{4c}-2^{4c})=v_5(149^4-16)+v_5(c)=1+v_5(c)$ meaning that we have that by LTE, $5^4 | c$ and $4 \cdot 5^4$ divides $n$. 

 Since $3^2$, $7^5$ and $4\cdot 5^4$ all divide $n$, the smallest value of $n$ working is their LCM, also $3^2 \cdot 7^5 \cdot 4 \cdot 5^4 = 2^2 \cdot 3^2 \cdot 5^4 \cdot 7^5$. Thus the number of divisors is $(2+1)(2+1)(4+1)(5+1) = \boxed{270}$.

 clarified by another user 

 notation note from another user

 \textbf{Note} We were able to use LTE with 3 and 7 but not 5 because in order to use LTE, we need \( p \mid x-y \).

 Obviously, \( 149^n \equiv 2^n \pmod{3} \) implies \( 149^n - 2^n \equiv 0 \pmod{3} \), so LTE works here.

 Furthermore, \( 149^n \equiv 2^n \pmod{7} \) implies \( 149^n - 2^n \equiv 0 \pmod{7} \), so LTE works here.

 However, when we get to the case of 5, we see that \( 149^n \equiv 2^n \pmod{5} \) doesn't always hold; specifically, this is only valid when \( n \) is a multiple of \( 4 \), which is why we let \( n = 4c \) in the solution.

 \textbf{mathboy282}

 





\textbf{Solution 1.2 (LTE + Binomial)}

Follow solution 1 to get that $3^2 \cdot 7^5$ is a divisor of $n$. We now take care of the $5^5$ part.Rewrite $149^n-2^n$ as $(150-1)^n-2^n$ and assume $n$ is even. Since this expression is divisible by $5^5$, we expand and get 
 \[\binom n0(1^n)-\binom n1(1^{n-1}150^1)+\binom n2(1^{n-2}150^2)\cdot\cdot\cdot+\binom nn(150^n)-2^n \equiv 0 \pmod {5^5}.\]Everything except the first three terms and the $2^n$ term is divisible by $5^5$, so we can rewrite this as
 \[1-150n+\frac{n(n-1)}{2}150^2-2^n \equiv 0 \pmod {5^5},\]And the following must be true:
 \[1-150n+\frac{n(n-1)}{2}150^2-2^n \equiv 0 \pmod {5^2}.\]Since $150n+\frac{n(n-1)}{2}150^2$ divides $5^2$, we can now see that
 \[1-2^n \equiv 0 \pmod {5^2}.\]Solving for $n$, $n \equiv 0 \pmod {2^2 \cdot 5}$.We also know this is true:
 \[1-150n+\frac{n(n-1)}{2}150^2-2^n \equiv 0 \pmod {5^3}.\]We know $n$ is divisible by $5$, so $150n+\frac{n(n-1)}{2}150^2$ divides $5^3$.And now,
 \[1-2^n \equiv 0 \pmod {5^3}\]gives 
 \[n \equiv 0 \pmod {2^2 \cdot 5^2}.\]This process can be repeated until we obtain $n \equiv 0 \pmod {2^2 \cdot 5^4}$. The desired solution is then simply $\text{lcm} (3^2 \cdot 7^5, 2^2 \cdot 5^4)$ which yields an answer of $(2+1)(5+1)(2+1)(4+1)=\boxed{270}$.

 





\textbf{Solution 2 (Simpler, just basic mods and Fermat's theorem)}

Note that for all $n$, $149^n - 2^n$ is divisible by $149-2 = 147$ by difference of $n$th powers. That is $3\cdot7^2$, so now we can clearly see that the smallest $n$ to make the expression divisible by $3^3$ is just $3^2$. Similarly, we can reason that the smallest $n$ to make the expression divisible by $7^7$ is just $7^5$. 

 Finally, for $5^5$, take $\pmod {5}$ and $\pmod {25}$ of each quantity (They happen to both be $-1$ and $2$ respectively, so you only need to compute once). One knows from \href{https://artofproblemsolving.com/wiki/index.php?title=Fermat\%27s\_theorem}{Fermat's theorem} that the maximum possible minimum $n$ for divisibility by $5$ is $4$, and other values are factors of $4$. Testing all of them(just $1$,$2$,$4$ using mods-not too bad), $4$ is indeed the smallest value to make the expression divisible by $5$, and this clearly is NOT divisible by $25$.Therefore, the smallest $n$ to make this expression divisible by $5^5$ is $2^2 \cdot 5^4$.

 Calculating the LCM of all these, one gets $2^2 \cdot 3^2 \cdot 5^4 \cdot 7^5$. Using the factor counting formula,the answer is $3\cdot3\cdot5\cdot6$ = $\boxed{270}$.

 





\textbf{Solution 3 (Using Totient Theorem for upper bound)}

For the $5^5$ case, we can see that the problem implies $149^n \equiv 2^n \pmod{5^5}.$ 

 The modular inverse of $2 \pmod{3125}$ is the least positive integer $x$ such that $2x \equiv 1 \pmod{3125}.$ 
 \[x = \frac{3125 + 1}{2} = 1563 \implies (149\cdot1563)^n \equiv 1 \pmod{5^5}.\]

 Now we can use Euler's Totient Theorem to get an upper bound for n, which states that if $\gcd(a,b) = 1,$ then $a^{\phi(b)} \equiv 1 \pmod{b}.$ $\phi(5^5) = 3125 \cdot (1 - \frac{1}{5}) = 2500,$ so $(149\cdot1563)^{2500} \equiv (1137)^{2500} \equiv 1 \pmod{5^5}.$

 For the number to end with $1,$ $n$ also needs to be a multiple of $4.$ So it's a factor of $2500$ that's also divisible by 4. We bash out the factors of $625.$ Through repeated squaring we see none of the factors work, so $2500$ must be the minimum. Therefore, n must be a multiple of $2500.$

 Next, we consider modulo $27$:
 \[149^n - 2^n \equiv 14^n - 2^n \pmod{27} \equiv 0 \pmod{27}.\]
 \[14^n \equiv 2^n \pmod{27}\] 
 \[7^n \equiv 1 \pmod{27}\] Apply Euler's Totient Theorem again to get $7^{18} \equiv 1 \pmod{27}.$ 

 Bashing the factors of 18, we find $18$ is the smallest value that works, so $n$ needs to be a multiple of $18.$

 For the $7^7$ case we can set up another congruence:
 \[149^n - 2^n \equiv 0 \pmod{7^7} \implies (2 + 7 \cdot 21)^n - 2^n \equiv 0 \pmod{7^7}\]

 To continue evaluating we use the Binomial Theorem:
 \[\sum_{k=0}^{n} \binom{n}{k} 2^{n-k} (7 \cdot 21)^k - 2^n \equiv 0 \pmod{7^7}\]
 \[\cancel{2^n} + \binom{n}{1} 2^{n-1} (7 \cdot 21) + \binom{n}{2} 2^{n-2} (7 \cdot 21)^2 + \dots \cancel{- 2^n} \equiv 0 \pmod{7^7}\]
 \[\binom{n}{1} \cdot 2^{n-1} \cdot 147 + \binom{n}{2} 2^{n-2} \cdot 147^2 + \binom{n}{3} 2^{n-3} \cdot 147^3 + \binom{n}{4} 2^{n-4} \cdot 147^4 \dots \equiv 0 \pmod{7^7}.\]

 Since $147 = 3 \cdot 7^2,$ the first term is a multiple of $7^2,$ the second is a multiple of $7^4,$ and the third is a multiple of $7^6.$ Everything after is a multiple of $7^7.$

 So $n$ must be a multiple of $7^5$ at the minimum for these three terms to also be multiples of $7^7.$

 Putting everything together we find the smallest value of $n$ to be $n = 2^2 \cdot 3^2 \cdot 5^4 \cdot 7^5,$ so the answer is $(2+1)(2+1)(4+1)(5+1) = \boxed{270}.$

 





\textbf{Solution 4 (Elementary and Thorough)}

As usual, denote $v_p(n)$ the highest power of prime $p$ that divides $n$. For divisibility by $3^3$, notice that $v_3(149^3 - 2^3) = 2$ as $149^3 - 2^3 =$ $(147)(149^2 + 2\cdot149 + 2^2)$, and upon checking mods, $149^2 + 2\cdot149 + 2^2$ is divisible by $3$ but not $9$. In addition, $149^9 - 2^9$ is divisible by $3^3$ because $149^9 - 2^9 = (149^3 - 2^3)(149^6 + 149^3\cdot2^3 + 2^6)$, and the rightmost factor equates to $1 + 1 + 1 \pmod{3} \equiv 0 \pmod{3}$. In fact, $n = 9 = 3^2$ is the least possible choice to ensure divisibility by $3^3$ because if $n = a \cdot 3^b$, with $3 \nmid a$ and $b < 2$, we write 
 \[149^{a \cdot 3^b} - 2^{a \cdot 3^b} =  (149^{3^b} - 2^{3^b})(149^{3^b(a - 1)} + 149^{3^b(a - 2)}\cdot2^{3^b}+\cdots2^{3^b(a - 1)}).\] Then, the rightmost factor is equivalent to $\pm a \pmod{3} \not\equiv 0 \pmod{3}$, and $v_3(149^{3^b} - 2^{3^b}) = b + 1 < 3$.

 For divisibility by $7^7$, we'll induct, claiming that $v_7(149^{7^k} - 2^{7^k}) = k + 2$ for whole numbers $k$. The base case is clear. Then, 
 \[v_7(149^{7^{k+1}} - 2^{7^{k+1}}) = v_7(149^{7^k} - 2^{7^k}) + v_7(149^{6\cdot7^k} + 2^{7^k}\cdot149^{5\cdot7^k} + \cdots + 2^{5\cdot7^k}\cdot149^{7^k} + 2^{6\cdot7^k}).\] By the induction hypothesis, $v_7(149^{7^k} - 2^{7^k}) = k + 2$. Then, notice that 
 \[S(k) = 149^{6\cdot7^k} + 2^{7^k}\cdot149^{5\cdot7^k} + \cdots + 2^{5\cdot7^k}\cdot149^{7^k} + 2^{6\cdot7^k} \equiv 7 \cdot 2^{6\cdot7^k}\pmod{7} \equiv 7 \cdot 2^{6\cdot7^k}\pmod{49}.\] This tells us that $S(k)$ is divisible by $7$, but not $49$ so that $v_7\left(S(k)\right) = 1$, completing our induction. We can verify that $7^5$ is the least choice of $n$ to ensure divisibility by $7^7$ by arguing similarly to the $3^3$ case. 

 Finally, for $5^5$, we take the powers of $149$ and $2$ in mod $5$ and mod $25$. Writing out these mods, we have that $149^n \equiv 2^n \pmod{5}$ if and only if $4 | n$, in which $149^n \equiv 2^n \equiv 1 \pmod{5}$. So here we claim that $v_5(149^{4\cdot5^k} - 2^{4\cdot5^k}) = k + 1$ and perform yet another induction. The base case is true: $5 | 149^4 - 2^4$, but $149^4 - 2^4 \equiv 1 - 16 \pmod{25}$. Now then, assuming the induction statement to hold for some $k$, 
 \[v_5(149^{4\cdot5^{k+1}} - 2^{4\cdot5^{k+1}}) = (k+1) + v_5(149^{4\cdot4\cdot5^k}+2^{4\cdot5^k}\cdot149^{3\cdot4\cdot5^k}+\cdots+2^{3\cdot4\cdot5^k}\cdot149^{4\cdot5^k}+2^{4\cdot4\cdot5^k}).\] Note that $S'(k) =  149^{4\cdot4\cdot5^k}+2^{4\cdot5^k}\cdot149^{3\cdot4\cdot5^k}+\cdots+2^{3\cdot4\cdot5^k}\cdot149^{4\cdot5^k}+2^{4\cdot4\cdot5^k}$ equates to $S''(k) = 1 + 2^{4\cdot5^k} + \cdots + 2^{16\cdot5^k}$ in both mod $5$ and mod $25$. We notice that $S''(k) \equiv 0 \pmod{5}$. Writing out the powers of $2$ mod $25$, we have $S''(0) \equiv 5 \pmod{25}$. Also $2^n \equiv 1 \pmod{25}$ when $n$ is a multiple of $20$. Hence for $k > 0$, $S''(k) \equiv 5 \mod{25}$. Thus, $v_5\left(S'(k)\right) = 1$, completing our induction. Applying the same argument from the previous two cases, $4\cdot5^4$ is the least choice to ensure divisibility by $5^5$.

 Our answer is the number of divisors of $\text{lcm}(3^2, 7^5, 2^2\cdot5^4) = 2^2 \cdot 3^2 \cdot 5^4 \cdot 7^5$. It is $(2 + 1)(2 + 1)(4 + 1)(5 + 1) = \boxed{270}$.

 





\textbf{Solution 5 (Official MAA)}

Analyze each prime power separately.Start with the case of $3^3$. By the Binomial Theorem,
 \begin{align*} 		149^n - 2^n &= (147+2)^n - 2^n \\ 		&= \binom n1 \cdot 147 \cdot 2^{n-1} 		+ \binom n2 \cdot 147^2 \cdot 2^{n-2}\\ 		&\qquad+ \binom n3 \cdot 147^3 \cdot 2^{n-3} 		+ \cdots. 	\end{align*}
Because $147$ is divisible by $3$, all terms after the first two are divisible by $3^3$, and the exponent of $3$ in the first term is less than that in the second term. Hence it is necessary and sufficient that $3^3 \mid 147n$, that is, $3^2 \mid n$.For the $7^7$ case, consider the same expansion as in the previous case. Because $147$ is divisible by $49 = 7^2$, all terms after the first three are divisible by $7^7$, and the exponent of $7$ in the first term is less than that in the second and third term. Hence it is necessary and sufficient that $7^7 \mid 147n$, that is, $7^5 \mid n$.For the $5^5$ case, working modulo $5$ gives $149^n - 2^n \equiv 4^n - 2^n = 2^n(2^n-1) \pmod 5$, so it must be that $4 \mid n$. Let $n = 4m$, and let $c = 149^4 - 2^4 = (149^2-2^2)(149^2+2^2) = 147 \cdot 151 \cdot 22205$. Note that $\frac c5$ is an integer not divisible by $5$. Expand by the Binomial Theorem again to get
 \begin{align*} 		(149^4)^m - (2^4)^m &= (c+16)^m - (16)^m \\ 		&= \binom m1 \cdot c \cdot 16^{m-1} 		+ \binom m2 \cdot c^2 \cdot 16^{m-2} \\ 		&\qquad+ \binom m3 \cdot c^3 \cdot 16^{m-3} 		+ \binom m4 \cdot c^4 \cdot 16^{m-4} 		+ \cdots. 	\end{align*}
All terms after the first four are divisible by $5^5$, and the exponent of $5$ in the first term is less than that in the second, third, or fourth term. Hence it is necessary and sufficient that $5^5 \mid cm$. Thus $5^4 \mid m$, and it follows that $4 \cdot 5^4 \mid n$.Therefore the least $n$ is $3^2 \cdot (2^2 \cdot 5^4) \cdot 7^5$. The requested number of divisors is $(1+2)(1+2)(1+4)(1+5) = 270$.

 The results of the above cases can be generalized using the following lemma.

 Lifting the Exponent Lemma: Let $p$ be an odd prime, and let $a$ and $b$ be integers relatively prime to $p$ such that $p \mid (a-b)$. Let $n$ be a positive integer. Then the number of factors of $p$ that divide $a^n - b^n$ is equal to the number of factors of $p$ that divide $a-b$ plus the number of factors of $p$ that divide $n$.

 




	\item (\textit{AIME I 2020}) Point $D$ lies on side $\overline{BC}$ of $\triangle ABC$ so that $\overline{AD}$ bisects $\angle BAC.$ The perpendicular bisector of $\overline{AD}$ intersects the bisectors of $\angle ABC$ and $\angle ACB$ in points $E$ and $F,$ respectively. Given that $AB=4,BC=5,$ and $CA=6,$ the area of $\triangle AEF$ can be written as $\frac{m\sqrt{n}}p,$ where $m$ and $p$ are relatively prime positive integers, and $n$ is a positive integer not divisible by the square of any prime. Find $m+n+p$.


 
 \textbf{Answer: }036



\textbf{Solution 1}

Points are defined as shown (Actually Point G is a point on line BC such that $\triangle AGD$ has circumcenter F, while Point H is a point on line BC such that $\triangle AHD$ has circumcenter E). Since $\angle{AFE} = \angle{AGD}$, and $\angle{AEF} = \angle{AHD}$, it is a breakthrough to get $\triangle AFE \sim \triangle AGH$. We'd like to compare altitudes to help us compare their areas. Now, note that $AD/2$ is the altitude of $\triangle AFE$, and Stewart's theorem implies that $AD/2 = \frac{\sqrt{18}}{2}$. Similarly, the altitude of $\triangle AGH$ is the altitude of $\triangle ABC$, or $\frac{3\sqrt{7}}{2}$. However, by the symmetry axis $BE$ on quadrilateral $ABHE$, we have $BH = AB = 4$, and similarly $CG = CA = 6$, so $GB = HC = 1$, and therefore $[AGH] = [ABC]$. From here, we get that the area of $\triangle AFE$ is $\frac{15\sqrt{7}}{4}*(\frac{\frac{\sqrt{18}}{2}}{\frac{3\sqrt{7}}{2} })^2=\frac{15\sqrt{7}}{14} \implies \boxed{036}$, by similarity. ~awang11 ~gougutheorem

 





\textbf{Solution 2}

Let $M_A$, $M_B$, $M_C$ be the midpoints of arcs $BC$, $CA$, $AB$. By Fact 5, we know that $M_AB = M_AC = M_AI$, and so by Ptolmey's theorem, we deduce that 
 \[AB\cdot M_AC + AC\cdot M_AB = BC\cdot M_AA \implies M_AA = 2M_AI.\]In particular, we have $AI = IM_A$. 
                    \begin{figure}[H]
                        \centering
                        \includegraphics[width=0.5\textwidth]{Solution/AIME_2020_I_13_Solution_Figure_1.png}
                    \end{figure}
                    Now the key claim is that: 

 \textbf{Claim:} $\triangle DEF$ and $\triangle M_AM_BM_C$ are homothetic at $I$ with ratio $2$. 

 \textit{Proof.} First, we show that $D$ is the midpoint of $M_AI$. Indeed, we have 
 \[\frac{ID}{DM_A} = \frac{BI}{BM_A}\cdot \frac{\sin\angle IBC}{\sin \angle CBM_A} = \frac{BI}{AI}\cdot\frac{\sin \angle B/2}{\sin \angle A/2} = 1\] by Ratio lemma and Law of Sines. 

 Now observe that:

 
                \begin{itemize}
                    \item $\overline{M_BM_C}$ is the perpendicular bisector of $\overline{AI}$,\item $\overline{EF}$ is the perpendicular bisector of $\overline{AD}$, and\item $ID = AI/2$.
                \end{itemize}


                Combining these facts gives that $\overline{EF}$ is a midline in $\triangle IM_BM_C$, which proves the claim. $\blacksquare$

 To finish, we compute $[M_AM_BM_C]$, noting that $[AEF] = [DEF] = \frac{1}{4}[M_AM_BM_C]$.

 By Heron's, we can calculate the circumradius $R = 8/\sqrt{7}$, and by Law of Cosines, we get 
 \begin{align*}\cos A &= \frac{9}{16}\implies \cos A/2 = \frac{5}{\sqrt{32}} \\ \cos B &= \frac{1}{8} \implies \cos B/2 = \frac{3}{4} \\ \cos C &= \frac{3}{4} \implies \cos C/2 = \sqrt{\frac{7}{8}}.\end{align*}
Then using $[XYZ] = 2R^2\sin X\sin Y\sin Z$, we can compute 
 \[[M_AM_BM_C] = 2\cdot \frac{64}{7}\cdot \frac{5}{\sqrt{32}}\cdot \frac{3}{4}\cdot \frac{\sqrt{7}}{\sqrt{8}} = \frac{30\sqrt{7}}{7}.\]Thus $[AEF] = 15\sqrt{7}/14$, which gives a final answer of $\boxed{036}$.

 





\textbf{Solution 3(coord bash + basic geometry)}

Let $\overline{BC}$ lie on the x-axis and $B$ be the origin. $C$ is $(5,0)$. Use Heron's formula to compute the area of triangle $ABC$. We have $s=\frac{15}{2}$. and $[ABC]=\sqrt{\frac{15 \cdot 7 \cdot 5 \cdot 3}{2^4}}=\frac{15\sqrt{7}}{4}$. We now find the altitude, which is $\frac{\frac{15\sqrt{7}}{2}}{5}=\frac{3\sqrt{7}}{2}$, which is the y-coordinate of $A$. We now find the x-coordinate of $A$, which satisfies $x^2 + (\frac{3\sqrt{7}}{2})^{2}=16$, which gives $x=\frac{1}{2}$ since the triangle is acute. Now using the Angle Bisector Theorem, we have $\frac{4}{6}=\frac{BD}{CD}$ and $BD+CD=5$ to get $BD=2$. The coordinates of D are $(2,0)$.Since we want the area of triangle $AEF$, we will find equations for perpendicular bisector of AD, and the other two angle bisectors. The perpendicular bisector is not too challenging: the midpoint of AD is $(\frac{5}{4}, \frac{3\sqrt{7}}{4})$ and the slope of AD is $-\sqrt{7}$. The slope of the perpendicular bisector is $\frac{1}{\sqrt{7}}$. The equation is(in point slope form) $y-\frac{3\sqrt{7}}{4}=\frac{1}{\sqrt{7}}(x-\frac{5}{4})$.The slope of AB, or in trig words, the tangent of $\angle ABC$ is $3\sqrt{7}$.Finding $\sin{\angle ABC}=\frac{\frac{3\sqrt{7}}{2}}{4}=\frac{3\sqrt{7}}{8}$ and $\cos{\angle ABC}=\frac{\frac{1}{2}}{4}=\frac{1}{8}$. Plugging this in to half angle tangent, it gives $\frac{\frac{3\sqrt{7}}{8}}{1+\frac{1}{8}}=\frac{\sqrt{7}}{3}$ as the slope of the angle bisector, since it passes through $B$, the equation is $y=\frac{\sqrt{7}}{3}x$.Similarly, the equation for the angle bisector of $C$ will be $y=-\frac{1}{\sqrt{7}}(x-5)$.For $E$ use the B-angle bisector and the perpendicular bisector of AD equations to intersect at $(3,\sqrt{7})$.For $F$ use the C-angle bisector and the perpendicular bisector of AD equations to intersect at $(\frac{1}{2}, \frac{9}{2\sqrt{7}})$.The area of AEF is equal to $\frac{EF \cdot \frac{AD}{2}}{2}$ since AD is the altitude of that triangle with EF as the base, with $\frac{AD}{2}$ being the height. $EF=\frac{5\sqrt{2}}{\sqrt{7}}$ and $AD=3\sqrt{2}$, so $[AEF]=\frac{15}{2\sqrt{7}}=\frac{15\sqrt{7}}{14}$ which gives $\boxed{036}$. NEVER overlook coordinate bash in combination with beginner synthetic techniques.~vvluo

 





\textbf{Solution 4 (Coordinate Bash + Trig)}


                    \begin{figure}[H]
                        \centering
                        \includegraphics[width=0.5\textwidth]{Solution/AIME_2020_I_13_Solution_Figure_1.png}
                    \end{figure}
                    

 Let $B=(0,0)$ and $BC$ be the line $y=0$.We compute that $\cos{\angle{ABC}}=\frac{1}{8}$, so $\tan{\angle{ABC}}=3\sqrt{7}$.Thus, $A$ lies on the line $y=3x\sqrt{7}$. The length of $AB$ at a point $x$ is $8x$, so $x=\frac{1}{2}$.

 We now have the coordinates $A=\left(\frac{1}{2},\frac{3\sqrt{7}}{2}\right)$, $B=(0,0)$ and $C=(5,0)$.We also have $D=(2,0)$ by the angle-bisector theorem and $M=\left(\frac{5}{4},\frac{3\sqrt{7}}{4}\right)$ by taking the midpoint.We have that because $\cos{\angle{ABC}}=\frac{1}{8}$, $\cos{\frac{\angle{ABC}}{2}}=\frac{3}{4}$ by half angle formula.

 We also compute $\cos{\angle{ACB}}=\frac{3}{4}$, so $\cos{\frac{\angle{ACB}}{2}}=\frac{\sqrt{14}}{4}$.

 Now, $AD$ has slope $-\frac{\frac{3\sqrt{7}}{2}}{2-\frac{1}{2}}=-\sqrt{7}$, so it's perpendicular bisector has slope $\frac{\sqrt{7}}{7}$ and goes through $\left(\frac{5}{4},\frac{3\sqrt{7}}{4}\right)$.

 We find that this line has equation $y=\frac{\sqrt{7}}{7}x+\frac{4\sqrt{7}}{7}$.

 As $\cos{\angle{CBI}}=\frac{3}{4}$, we have that line $BI$ has form $y=\frac{\sqrt{7}}{3}x$.Solving for the intersection point of these two lines, we get $x=3$ and thus $E=\left(3, \sqrt{7}\right)$

 We also have that because $\cos{\angle{ICB}}=\frac{\sqrt{14}}{4}$, $CI$ has form $y=-\frac{x\sqrt{7}}{7}+\frac{5\sqrt{7}}{7}$.

 Intersecting the line $CI$ and the perpendicular bisector of $AD$ yields $-\frac{x\sqrt{7}}{7}+\frac{5\sqrt{7}}{7}=\frac{x\sqrt{7}}{7}+\frac{4\sqrt{7}}{7}$.

 Solving this, we get $x=\frac{1}{2}$ and so $F=\left(\frac{1}{2},\frac{9\sqrt{7}}{14}\right)$.

 We now compute $EF=\sqrt{\left(\frac{5}{2}\right)^2+\left(\frac{5\sqrt{7}}{14}\right)^2}=\frac{5\sqrt{14}}{7}$.We also have $MA=\sqrt{\left(\frac{3}{4}\right)^2+\left(\frac{3\sqrt{7}}{4}\right)^2}=\frac{3\sqrt{2}}{2}$.

 As ${MA}\perp{EF}$, we have $[\triangle{AEF}]=\frac{1}{2}\left(\frac{3\sqrt{2}}{2}\times\frac{5\sqrt{14}}{7}\right)=\frac{15\sqrt{7}}{14}$.

 
The desired answer is $15+7+14=\boxed{036}$ ~Imayormaynotknowcalculus

 





\textbf{Solution 5 (Barycentric Coordinates)}


                    \begin{figure}[H]
                        \centering
                        \includegraphics[width=0.5\textwidth]{Solution/AIME_2020_I_13_Solution_Figure_1.png}
                    \end{figure}
                    

 As usual, we will use homogenized barycentric coordinates.

 We have that $AD$ will have form $3z=2y$. Similarly, $CF$ has form $5y=6x$ and $BE$ has form $5z=4x$.Since $A=(1,0,0)$ and $D=\left(0,\frac{3}{5},\frac{2}{5}\right)$, we also have $M=\left(\frac{1}{2},\frac{3}{10},\frac{1}{5}\right)$.It remains to determine the equation of the line formed by the perpendicular bisector of $AD$.

 This can be found using EFFT. Let a point $T$ on $EF$ have coordinates $(x, y, z)$.We then have that the displacement vector $\overrightarrow{AD}=\left(-1, \frac{3}{5}, \frac{2}{5}\right)$ and that the displacement vector $\overrightarrow{TM}$ has form $\left(x-\frac{1}{2},y-\frac{3}{10},z-\frac{1}{5}\right)$.Now, by EFFT, we have $5^2\left(\frac{3}{5}\times\left(z-\frac{1}{5}\right)+\frac{2}{5}\times\left(y-\frac{3}{10}\right)\right)+6^2\left(-1\times\left(z-\frac{1}{5}\right)+\frac{2}{5}\times\left(x-\frac{1}{2}\right)\right)+4^2\left(-1\times\left(y-\frac{3}{10}\right)+\frac{3}{5}\times\left(x-\frac{1}{2}\right)\right)=0$.This equates to $8x-2y-7z=2$.

 Now, intersecting this with $BE$, we have $5z=4x$, $8x-2y-7z=2$, and $x+y+z=1$.This yields $x=\frac{2}{3}$, $y=-\frac{1}{5}$, and $z=\frac{8}{15}$, or $E=\left(\frac{2}{3},-\frac{1}{5},\frac{8}{15}\right)$.

 Similarly, intersecting this with $CF$, we have $5y=6x$, $8x-2y-7z=2$, and $x+y+z=1$.Solving this, we obtain $x=\frac{3}{7}$, $y=\frac{18}{35}$, and $z=\frac{2}{35}$, or $F=\left(\frac{3}{7},\frac{18}{35},\frac{2}{35}\right)$.

 We finish by invoking the Barycentric Distance Formula twice; our first displacement vector being $\overrightarrow{FE}=\left(\frac{5}{21},-\frac{5}{7},\frac{10}{21}\right)$.We then have $FE^2=-25\left(-\frac{5}{7}\cdot\frac{10}{21}\right)-36\left(\frac{5}{21}\cdot\frac{10}{21}\right)-16\left(\frac{5}{21}\cdot-\frac{5}{7}\right)=\frac{50}{7}$, thus $FE=\frac{5\sqrt{14}}{7}$.

 Our second displacement vector is $\overrightarrow{AM}=\left(-\frac{1}{2},\frac{3}{10},\frac{1}{5}\right)$.As a result, $AM^2=-25\left(\frac{3}{10}\cdot\frac{1}{5}\right)-36\left(-\frac{1}{2}\cdot\frac{1}{5}\right)-16\left(-\frac{1}{2}\cdot\frac{3}{10}\right)=\frac{9}{2}$, so $AM=\frac{3\sqrt{2}}{2}$.

 As ${AM}\perp{EF}$, the desired area is $\frac{\frac{5\sqrt{14}}{7}\times\frac{3\sqrt{2}}{2}}{2}={\frac{15\sqrt{7}}{14}}\implies{m+n+p=\boxed{036}}$. ~Imayormaynotknowcalculus

 
\textbf{Remark}: The area of $\triangle{AEF}$ can also be computed using the \textit{Barycentric Area Formula}, although it may increase the risk of computational errors; there are also many different ways to proceed once the coordinates are determined.

 





\textbf{Solution 6 (geometry+trig)}


                    \begin{figure}[H]
                        \centering
                        \includegraphics[width=0.5\textwidth]{Solution/AIME_2020_I_13_Solution_Figure_1.png}
                    \end{figure}
                    

 To get the area of $\triangle AEF$, we try to find $AM$ and $\angle EAF$.

 Since $AD$ is the angle bisector, we can get that $BD=2$ and $CD=3$. By applying Stewart's Theorem, we can get that $AD=3\sqrt{2}$. Therefore $AM=\frac{3\sqrt{2}}{2}$.

 Since $EF$ is the perpendicular bisector of $AD$, we know that $AE = DE$. Since $BE$ is the angle bisector of $\angle BAC$,we know that $\angle ABE = \angle DBE$. By applying the Law of Sines to $\triangle ABE$ and $\triangle DBE$, we know that$\sin \angle BAE = \sin \angle BDE$.  Since $BD$ is not equal to $AB$ and therefore these two triangles are not congruent, we know that $\angle BAE$ and $\angle BDE$ are supplementary.  Then we know that $\angle ABD$ and $\angle AED$ are also supplementary. Given that $AE=DE$, we can get that $\angle DAE$ is half of $\angle ABC$.  Similarly, we have $\angle DAF$ is half of $\angle ACB$.

 By applying the Law of Cosines, we get $\cos \angle ABC = \frac{1}{8}$, and then $\sin \angle ABC = \frac{3\sqrt{7}}{8}$.  Similarly, we can get $\cos \angle ACB = \frac{3}{4}$ and $\sin \angle ACB = \frac{\sqrt{7}}{4}$. Based on some trig identities, we can compute that $\tan \angle DAE = \frac{\sin \angle ABC}{1 + \cos \angle ABC} = \frac{\sqrt{7}}{3}$, and $\tan \angle DAF = \frac{\sqrt{7}}{7}$.

 Finally, the area of $\triangle AEF$ equals $\frac{1}{2}AM^2(\tan \angle DAE + \tan \angle DAF)=\frac{15\sqrt{7}}{14}$. Therefore, the final answer is $15+7+14=\boxed{036}$. ~xamydad

 \textbf{Remark}: I didn't figure out how to add segments $AF$, $AE$, $DF$ and $DE$. Can someone please help add these segments? 

 (Added :) ~Math\_Genius\_164)。

 





\textbf{Solution 7}


                    \begin{figure}[H]
                        \centering
                        \includegraphics[width=0.5\textwidth]{Solution/AIME_2020_I_13_Solution_Figure_1.png}
                    \end{figure}
                    First and foremost $\big[\triangle{AEF}\big]=\big[\triangle{DEF}\big]$ as $EF$ is the perpendicular bisector of $AD$. Now   note that quadrilateral $ABDF$ is cyclic, because $\angle{ABF}=\angle{FBD}$ and $FA=FD$. Similarly quadrilateral $AEDC$ is cyclic, 
 \[\implies \angle{EDA}=\dfrac{C}{2}, \quad \angle{FDA}=\dfrac{B}{2}\]Let $A'$,$B'$, $C'$ be the $A$,$B$, and $C$ excenters of $\triangle{ABC}$ respectively. Then it follows that $\triangle{DEF} \sim \triangle{A'C'B'}$. By angle bisector theorem we have $BD=2 \implies \dfrac{ID}{IA}=\dfrac{BD}{BA}=\dfrac{1}{2}$. Now let the feet of the perpendiculars from $I$ and $A'$ to $BC$ be $X$ and $Y$ resptively. Then by tangents we have 
 \[BX=s-AC=\dfrac{3}{2} \implies XD=2-\dfrac{3}{2}=\dfrac{1}{2}\] 
 \[CY=s-AC \implies YD=3-\dfrac{3}{2}=\dfrac{3}{2} \implies \dfrac{ID}{DA'}=\dfrac{XD}{YD}=\dfrac{1}{3} \implies \big[\triangle{DEF}\big]=\dfrac{1}{16}\big[\triangle{A'C'B'}\big]\] From the previous ratios, $AI:ID:DA'=2:1:3 \implies AD=DA' \implies \big[\triangle{ABC}\big]=\big[\triangle{A'BC}\big]$ Similarly we can find that $\big[\triangle{B'AC}\big]=2\big[\triangle{ABC}\big]$ and $\big[\triangle{C'AB}\big]=\dfrac{4}{7}\big[\triangle{ABC}\big]$ and thus 
 \[\big[\triangle{A'B'C'}\big]=\bigg(1+1+2+\dfrac{4}{7}\bigg)\big[\triangle{ABC}\big]=\dfrac{32}{7}\big[\triangle{ABC}\big] \implies \big[\triangle{DEF}\big]=\dfrac{2}{7}\big[\triangle{ABC}\big]=\dfrac{15\sqrt{7}}{14} \implies m+n+p = \boxed{036}\]-tkhalid

 





\textbf{Solution 8 -Trigonometry(only)}


                    \begin{figure}[H]
                        \centering
                        \includegraphics[width=0.5\textwidth]{Solution/AIME_2020_I_13_Solution_Figure_1.png}
                    \end{figure}
                    Trig values we use here:

 $\cos A = \frac{9}{16}$

 $\cos \frac{A}{2} = \frac{5}{4\sqrt2}$

 $\sin \frac{A}{2} = \frac{\sqrt7}{4\sqrt2}$

 $\cos \frac{B}{2} = \frac{3}{4}$

 $\cos \frac{C}{2} = \frac{\sqrt7}{2\sqrt2}$

 
First let the incenter be $I$. Let $M$ be the midpoint of minor arc $BC$ on $(ABC)$ and let $K$ be the foot of $M$ to $BC$.

 We can find $AD$ using Stewart's Theorem: from Angle Bisector Theorem $BD = 2$ and $CD = 3$. Then it is easy to find that $AD = 3\sqrt3$.

 Now we trig bash for $DI = MI - MD$. Notice that $MI = MB$ from the Incenter Excenter Lemma. We obtain that $MB = \frac{BK}{\cos \frac{A}{2}} = \frac{\frac{5}{2}}{\frac{5}{4\sqrt2}}=2\sqrt2$. To get $MD$ we angle chase to get $\angle KDM = \frac{A}{2}+C$. Then 
 \[\cos(\frac{A}{2}+C) = \cos\frac{A}{2}\cos C - \sin\frac{A}{2}\sin C = \frac{1}{2\sqrt2} = \frac{\frac{1}{2}}{MD}\] gives $MD = \sqrt{2}$. This means $DI = \sqrt2$.

 Now let $AI \cap EF = G$. It is easy to angle chase $\angle GIE = 90- \frac{B}{2}$ and $\angle GIF = 90- \frac{C}{2}$. Since $GI = GD - ID = \frac{3\sqrt2}{2}-\sqrt2=\frac{\sqrt2}{2}$, we compute that 
 \[EF = EG + FG = \frac{\sqrt2}{2}(\cot B/2 + \cot C/2) = \frac{\sqrt2}{2}(\frac{3}{\sqrt7}+\sqrt7) = \frac{5\sqrt{14}}{7}\] which implies 
 \[[AEF] = AG*EF/2 = \frac{3\sqrt2}{2} * \frac{5\sqrt{14}}{7} / 2 = \frac{15\sqrt7}{14}\] which gives an answer of $\boxed{36}$. ~Leonard\_my\_dude~

 





\textbf{Solution 9 (Official MAA 1)}

Let $x = \angle BAD = \angle CAD$, $y = \angle CBE = \angle ABE$, and $z = \angle BCF = \angle ACF$. Notice that $x+y+z = 90^\circ$.

 In $\triangle ABD$, segment $\overline{BE}$ is the bisector of $\angle ABD$, and $E$ lies on the perpendicular bisector of side $\overline{AD}$. Therefore $E$ is the midpoint of arc $\stackrel{\textstyle\frown}{AD}$ on the circumcircle of $\triangle ABD$. It follows that $\angle BED = \angle BAD = x$ and $\angle EDA = \angle EBA = y$. Likewise, $ACDF$ is cyclic, $\angle CFD = \angle CAD = x$, and $\angle FDA = \angle FCA = z$. Because $\overline{EF}$ is the perpendicular bisector of $\overline{AD}$, triangles $AEF$ and $DEF$ are congruent, implying that
 \begin{align*} [\triangle AEF] = [\triangle DEF] &= \frac{DE\cdot DF\cdot\sin(\angle EDF)}{2}\\ &= \frac{DE\cdot DF\cdot\sin(y+z)}{2} = \frac{DE\cdot DF\cdot\cos x}{2}. \end{align*}

                    \begin{figure}[H]
                        \centering
                        \includegraphics[width=0.5\textwidth]{Solution/AIME_2020_I_13_Solution_Figure_1.png}
                    \end{figure}
                    

 Applying the Law of Sines to $\triangle BED$ and $\triangle CFD$ gives
 \[DE = BD\cdot\frac{\sin y}{\sin x}\text{~ and ~} DF = CD\cdot\frac{\sin z}{\sin x}.\]By the Angle Bisector Theorem, $BD = 2$ and $CD = 3$. Combining the above information yields
 \[[\triangle AEF] = \frac{3\sin y\cdot\sin z\cdot\cos x}{\sin^2 x}.\]Applying the Law of Cosines to $\triangle ABC$ gives $\cos 2x = \frac9{16}$, $\cos 2y = \frac1{8}$, and $\cos 2z = \frac34$. By the Half Angle Formulas,
 \[\sin^2x = \frac7{32},~~ \cos x = \sqrt{\frac{25}{32}},~~ \sin y = \sqrt{\frac7{16}}, \text{~ and ~} \sin z = \sqrt{\frac18}.\]Therefore
 \[[\triangle AEF] = \frac{3\cdot\sqrt{\frac{7}{16}}\cdot\sqrt{\frac{1}{8}}\cdot\sqrt{\frac{25}{32}}} {\frac{7}{32}} = \frac{15\sqrt{7}}{14}.\]The requested sum is $15+7+14 = 36$.

 





\textbf{Solution 10 (Official MAA 2)}

Let the point $M$ be the midpoint of $\overline{AD}$, let $I$ be the incenter of $\triangle ABC$ which is the common point of lines $AD$, $BE$, and $CF$, and let $r$ be the inradius of $\triangle ABC$. The semiperimeter of $\triangle ABC$ is
 \[s = \frac{AB + BC + CA}2 = \frac{15}2,\]and Heron's Formula gives the area of $\triangle ABC$ as
 \[\sqrt{s(s-AB)(s-BC)(s-CA)} = \frac{15\sqrt7}4.\]This area is also $rs$ implying that $r = \frac{\sqrt7}2$. Stewart's Theorem gives $AD =3\sqrt2$. Because the ratio of the areas of $\triangle IBC$ and $\triangle ABC$ is $\frac{ID}{AD},$ it follows that
 \[ID = AD\cdot\frac{\frac{r\cdot BC}2}{\frac{15\sqrt7}4} = \sqrt2.\]Thus $IM = MD - ID = \frac{\sqrt2}2$.

 
                    \begin{figure}[H]
                        \centering
                        \includegraphics[width=0.5\textwidth]{Solution/AIME_2020_I_13_Solution_Figure_1.png}
                    \end{figure}
                    

 Note that $\angle EFI = 90^{\circ} - \angle FIA = 90^{\circ} - \angle CID = 90^{\circ} - \frac{\angle A}{2} - \frac{\angle C}{2}  = \frac{\angle B}{2} = \angle IBC$. Thus$\triangle IBC \sim \triangle IFE$. The height of $\triangle IBC$ to $I$ is $r=\frac{\sqrt7}2$, and the height of $\triangle IFE$ to $I$ is $IM=\frac{\sqrt2}2$, so $EF = BC\cdot \frac{IM}r = \frac{5\sqrt{14}}{7}$. The needed area of $\triangle AEF$ is $\frac12\cdot EF\cdot \frac{AD}2 = \frac{15\sqrt7}{14}$, as above.

 





\textbf{Solution 11 Bash for life}

Firstly, it is easy to find $BD=2,CD=3$ with angle bisector theorem.

 Using LOC and some trig formulas we get all those values:$cos\angle{B}=\frac{1}{8},sin\angle{B}=\frac{3\sqrt{7}}{8},cos\angle{\frac{B}{2}}=\frac{3}{4},cos\angle{\frac{{ACB}}{2}}=\frac{\sqrt{14}}{4}$Now we find the coordinates of points $A,E,F$ and we apply shoelace theorem later. Point A's coordinates is $(4*\frac{1}{8},4*\frac{3\sqrt{7}}{8})=(\frac{1}{2},\frac{3\sqrt{7}}{2})$, Let $AJ$ is perpendicular to $BC$, $tan\angle{JAD}=\frac{2-\frac{1}{2}}{\frac{3\sqrt{7}}{2}}=\frac{\sqrt{7}}{7}$, which means the slope of $FE$ is $\frac{\sqrt{7}}{7}$. Find the coordinate of $M$, it is easy, $(\frac{5}{4},\frac{3\sqrt{7}}{4})$, the function $EF$ is $y=\frac{\sqrt{7}x}{7}+\frac{4\sqrt{7}}{7}$. Now find the intersection of $EF,EB$, $\frac{\sqrt{7}x}{7}+\frac{4\sqrt{7}}{7}=\frac{\sqrt{7}x}{3}$, getting that $x=3,E(3,\sqrt{7})$Now we look at line segment $CF$, since $cos\angle{\frac{ACB}{2}}=\frac{\sqrt{14}}{4},tan\angle{\frac{ACB}{2}}=\frac{\sqrt{7}}{7}$. Since the line passes $(5,0)$, we can set the equation to get the $CF:y=-\frac{\sqrt{7}}{7}+\frac{5\sqrt{7}}{7}$, find the intersection of $CF,EF$,$-\frac{\sqrt{7}x}{7}+\frac{5\sqrt{7}}{7}=\frac{\sqrt{7}x}{7}+\frac{4\sqrt{7}}{7}$, getting that $F:(\frac{1}{2},\frac{9\sqrt{7}}{14})$and in the end we use shoelace theorem with coordinates of $A,F,E$ getting the area $\frac{15\sqrt{7}}{14}$ leads to the final answer $\boxed{36}$

 





\textbf{Solution 12 (Simple geometry)}


                    \begin{figure}[H]
                        \centering
                        \includegraphics[width=0.5\textwidth]{Solution/AIME_2020_I_13_Solution_Figure_1.png}
                    \end{figure}
                    Using the <i>\textbf{Claim}</i>  (below) we get $I$ is orthocenter of $\triangle DEF,\angle EDG = \beta,$$\angle FDG = \gamma.$ So area of $\triangle DEF$ is
 \[[\triangle DEF] =\frac {DG \cdot FE}{2} = \frac {AD^2}{8} (\tan \beta + \tan \gamma).\]

 Semiperimeter of $\triangle ABC \hspace{10mm} s = \frac{15}{2},$ so the bisector
 \[AD= \frac{2}{AB + AC}\sqrt{s(s-BC) \cdot AB\cdot AC} = 3\sqrt{2}.\]

 We get the inradius by applying Heron's formula 
 \[r = \sqrt{\frac {(s-AB)(s-BC)(s-AC)}{s}} = \sqrt {\frac {3\cdot 5 \cdot 7}{15 \cdot 4}}=\frac {\sqrt {7}}{2}.\]

 We use formulas for inradius and get 
 \[\tan \beta = \frac{r}{s – AC} = \frac {\sqrt {7}}{3} , \hspace{10mm} \tan \gamma  = \frac{r}{s – AB} = \frac {\sqrt {7}}{7}.\]The area $\hspace{30mm} [\triangle DEF] = \frac{15 \sqrt 7}{14}.$

 
                    \begin{figure}[H]
                        \centering
                        \includegraphics[width=0.5\textwidth]{Solution/AIME_2020_I_13_Solution_Figure_1.png}
                    \end{figure}
                    <i>\textbf{Claim}</i> 

 Let $I$ be incenter of $\triangle ABC.$ Then bisector $\overline{BI},$ perpendicular bisector of $\overline{AD},$ and perpendicular dropped to bisector $\overline{CI}$ from point $D$ are concurrent.

 <i>\textbf{Proof}</i>

 Denote $\angle BAC = 2 \alpha,  \angle ABC = 2\beta,\angle ACB = 2 \gamma.$Then $\alpha + \beta + \gamma = 90^\circ.$Denote $P$  the intersection point of $BC$ and the tangent line to the circumcircle at point $A.$

 WLOC,$\hspace{20mm}\gamma > \beta$ (case  $\gamma = \beta$ is trivial).

 $\angle PAC = \angle ABC = 2 \beta$ (this angles are measured by half the arc $\overset{\Large\frown} {AC}$  of the circumcircle).

 
 \[\angle APC = \angle ACD - \angle PAC = 2(\gamma - \beta),\] 
 \[\angle ADP = 180^o – \angle DAC - \angle ACD = 180^o – \alpha – 2 \gamma = \alpha + 2 \beta = \angle DAP \implies AP = DP.\] Therefore bisector of angle P coincite with the perpendicular bisector of $\overline{AD}$.

 By applying the Law of Sines to $\triangle ABP$ we get $\hspace{20mm}\frac {BP}{AP} =  \frac {\sin 2 \gamma}{\sin 2 \beta.}$

 Let $E$ be crosspoint of $PG$  and bisector $BI.$ By applying the Law of Sines to $\triangle BEP$ we get 
 \[\frac {EP}{BP} = \frac {\sin \beta}{\sin(180^o – \beta – (\gamma – \beta))}= \frac {\sin \beta}{\sin \gamma}.\]

 Let $E'$ be crosspoint of $PG$ and the perpendicular dropped to bisector $\overline{CI}$ from point $D.$ 
 \[\frac {E'P}{DP} = \frac {\sin (90^\circ – \gamma)}{\sin(180^o - (90^o – \gamma) – (\gamma - \beta))}= \frac {\cos\gamma}{\cos \beta}.\]$\hspace{50mm}\frac {E'P} {EP} = \frac {E'P}{DP} \cdot \frac {AP}{BP} \cdot \frac {BP}{EP} = \frac {\cos\gamma}{\cos \beta} \cdot \frac {\sin 2\beta}{\sin 2\gamma}\cdot \frac {\sin\gamma}{\sin \beta} = 1 \implies  E$  coincide with $E'.$

 \textbf{vladimir.shelomovskii@gmail.com, vvsss}

 





\textbf{Solution 13 (Trig)}

Denote $I$ as the incenter of $\triangle ABC$ and $M$ as the intersecting point of $AD$ and $EF$. 

 We solve this question based on this equation: $[\triangle AEF]=\dfrac{1}{2}\times AM\times (\dfrac{IM}{\tan{\angle EFI}}+\dfrac{IM}{\tan{\angle FEI}})$.

 
Firstly, let's compute the trig value of $\angle EFI$ and $\angle FEI$. Using the Cosine Law and half-angle formula, we obtain:
 \[\cos{\angle B}=\dfrac{1}{8}, \sin{\dfrac{\angle B}{2}}=\dfrac{\sqrt{7}}{4}, \cos{\dfrac{\angle B}{2}}=\dfrac{3}{4}\]
 \[\cos{\angle C}=\dfrac{3}{4}, \sin{\dfrac{\angle C}{2}}=\dfrac{1}{\sqrt{8}}, \cos{\dfrac{\angle C}{2}}=\dfrac{\sqrt{7}}{\sqrt{8}}\]

 Noticing that by angle chasing, $\angle EFI= 90^{\circ}-\angle FIA=90^{\circ}-\angle CID=90^{\circ}-(180^{\circ}-\dfrac{\angle A}{2}-\angle B-\dfrac{\angle C}{2})= \dfrac{\angle B}{2}$. Similarly, $\angle FEI=\dfrac{\angle C}{2}$.Thus, $\tan{\angle EFI}=\dfrac{\sqrt{7}}{3}$ and $\tan{\angle FEI}=\dfrac{1}{\sqrt{7}}$

 Then we find $AM$: by Stewart's theorem, $AM=\dfrac{AD}{2}=\dfrac{3\sqrt{2}}{2}$.

 For $MI$, first noticing that $MI=MD-ID=\dfrac{3\sqrt{2}}{2}-ID$. Then, by Angel Bisector's theorem, we have $BD=2, CD=3$, and applying Sine Law: 
 \[\dfrac{ID}{\sin{\dfrac{\angle C}{2}}}=\dfrac{CD}{\sin{\angle CID}}\]
 \[\dfrac{ID}{\sin{\dfrac{\angle C}{2}}}=\dfrac{CD}{\cos{\dfrac{\angle B}{2}}}\]Substituting the known values, we get $ID=\sqrt{2}$. Thus, $MI=\dfrac{\sqrt{2}}{2}$

 Now, we already find all the unknown values in the very beginning equation. Substituting all the values, we obtain $[\triangle AEF]=\dfrac{15\sqrt{7}}{14}$, which gives a final answer of $\boxed{036}$.

 





\textbf{Solution 14 (No Trig, No Coordinate Bashing, Simple Geometry)}


                    \begin{figure}[H]
                        \centering
                        \includegraphics[width=0.5\textwidth]{Solution/AIME_2020_I_13_Solution_Figure_1.png}
                    \end{figure}
                    [AUTHOR'S NOTE: Someone please help make a better diagram...]

 By Angle Bisector Theorem, $BD=2,DC=3$. By Stewart's, $AD=3\sqrt2$, so $AK=\dfrac{3\sqrt2}{2}$. Thus, it suffices to find $FE$. 

 Notice that $\triangle AIK \cong \triangle DIK$, because $IK$ is the perpendicular bisector of $AD$. Also notice that $\triangle AIK \cong \triangle AJK$ by $SAS$. Thus, $AJDI$ is a rhombus, and $AI \parallel JD, AJ\parallel ID$.

 Now notice $\triangle CJD \sim \triangle CAB$ by $AAA$. Since $\dfrac{CD}{CB}=\dfrac{JD}{AB}=\dfrac{3}{5}$, we have that $JD=ID=IA=AJ=\dfrac{12}{5}$. Thus, by Pythagorean Theorem, $KJ=\dfrac{3\sqrt{14}}{10},IJ=\dfrac{3\sqrt{14}}{5}.$

 We now do some length chasing. $BI=4-\dfrac{12}{5}=\dfrac{8}{5},CJ=6-\dfrac{12}{5}=\dfrac{18}{5}.$ By Angle Bisector Theorem on $\triangle BID,$ we have $IL=\dfrac{16}{15}$. Similarly on $\triangle CJD$, $JM=\dfrac{72}{55}.$ We also have $AG=\dfrac{24}{11}$ and $AH=\dfrac{8}{3}$ by Angle Bisector Theorem on $\triangle ABC$. Thus, $GI=AI-AG=\dfrac{12}{5}-\dfrac{24}{11}=\dfrac{12}{55}$, and $JH=AH-AJ=\dfrac{8}{3}-\dfrac{12}{5}=\dfrac{4}{15}.$

 It is easily shown that $\triangle GIF \sim \triangle JMF$, so by the side length ratios, we have 
 \[\dfrac{JM}{GI}=\dfrac{JF}{IF}=\dfrac{\frac{72}{55}}{\frac{12}{55}}=6.\]Similarly, we can see that $\triangle EJH \sim \triangle EIL$, so by the side length ratios, we have
 \[\dfrac{IL}{JK}=\dfrac{EI}{JE}=\dfrac{\frac{16}{15}}{\frac{4}{15}}=4.\]So, we can see that the length of $FE$ is $\dfrac{6}{7}+\dfrac{1}{3}=\dfrac{25}{21}$ times the length of $IJ$. However, the length of $IJ$ is $\dfrac{3\sqrt{14}}{5}$, so $FE=\dfrac{5\sqrt{14}}{7}$, and the area of $\triangle AFE$ is $\dfrac{1}{2}\times \dfrac{5\sqrt{14}}{7}\times \dfrac{3\sqrt2}{2}=\dfrac{15\sqrt7}{14}.$ Thus, the desired answer is $15+7+14=\boxed{036}.$

 


 





\textbf{Solution 15(Similar to Solutions with analytic geometry but this is pure coordinate bash)}

Set $B$ to be at $(0, 0)$. Then $C$ will be at $(5, 0)$. Let $A$ be at $(x, y)$. Hence, 

 $x^{2} + y^{2} = 16$

 $(x - 5)^{2} + y^{2} = 36$

 Subtracting yields

 $x^{2} - 10x + 25 - x^{2} = 20 \implies 10x = 5 \implies x = \frac{1}{2}$

 Then $y = 16 - \frac{1}{4} = \frac{63}{4} \implies y = \frac{3\sqrt{7}}{2}$

 Now by the Angle Bisector Theorem, $\frac{AB}{AC} = \frac{BD}{DC} \implies \frac{2}{3} = \frac{x}{5 - x}$ since $BD + DC = 5$. Now $10 - 2x = 3x \implies x = 2$. Hence $D$ is at $(2, 0)$. Let the point where $BE$ intersects $AC$ be $G$ and likewise $H$ for the point where $CF$ intersects $AB$. By the Angle Bisector Theorem, 

 $\frac{AB}{BC} = \frac{AG}{GB} \implies \frac{4}{5} = \frac{x}{6 - x}$. Solving yields $x = \frac{8}{3}$ and so $AG = \frac{8}{3}$ and $GC = \frac{10}{3}$. Doing similar things to finding the length of $AH$, we get $AH = \frac{24}{11}$ and $HB = \frac{20}{11}$. 

 Now we find the equation of line $AC$. The slope is just $\frac{\frac{3\sqrt{7}}{2}}{\frac{1}{2} - 5} = -\frac{\sqrt{7}}{3}$. Hence the equation of $AC$ is $y = -\frac{x\sqrt{7}}{3} + b$ and we can plug in point $(5, 0)$ because it lies on line $AC$ to get $0 = -\frac{5\sqrt{7}}{3} + b$ and solve to get the equation of line $AC$ is just $y = -\frac{x\sqrt{7}}{3} + \frac{5\sqrt{7}}{3}$.

 Therefore, any point on line $AC$ must have the form $(x, -\frac{x\sqrt{7}}{3} + \frac{5\sqrt{7}}{3})$. Let $G$ be at point $(x, -\frac{x\sqrt{7}}{3} + \frac{5\sqrt{7}}{3})$. We know $GC$ is just $\frac{10}{3}$ so the distance between points $(x, -\frac{x\sqrt{7}}{3} + \frac{5\sqrt{7}}{3})$ and $(5, 0)$ is just $\frac{10}{3}$. Solving, we get $x = \frac{5}{2}, \frac{15}{2}$ but we know $x$ can at most be $5$ hence only $x = \frac{5}{2}$ works. Thus, $G$ will be at point $(\frac{5}{2}, \frac{5\sqrt{7}}{6})$. 

 The same process applies for finding the coordinates of point $H$. We end up with $H$ being at point $(\frac{5}{22}, \frac{15\sqrt{7}}{22})$.

 
We now find the equation of $BE$. Note that the slope is just $\frac{\frac{5\sqrt{7}}{6}}{\frac{5}{2}} = \frac{\sqrt{7}}{3}$. We know that $B$ is at the origin and so the equation of line $BE$ is just $y = \frac{x\sqrt{7}}{3}$. 

 Then find the equation of $CF$ in a similar way. We get $y = -\frac{x\sqrt{7}}{7} + \frac{5\sqrt{7}}{7}$.

 Finally, the last step is finding the equation of the perpendicular bisector of $AD$. Note that the slope of $AD$ is just $\frac{\frac{3\sqrt{7}}{2}}{\frac{1}{2} - 2} = -\sqrt{7}$. Hence the slope of the perpendicular bisector is just $\frac{\sqrt{7}}{7}$ because that is the negative reciprocal of $-\sqrt{7}$. We also know that because this is the bisector, it passes through the midpoint of $AD$ which is just $(\frac{\frac{1}{2} + 2}{2}, \frac{3\sqrt{7}}{4}) = (\frac{5}{4}, \frac{3\sqrt{7}}{4})$. Therefore $y = -x\sqrt{7} + b$ and we know $(\frac{5}{4}, \frac{3\sqrt{7}}{4})$ is on it and plugging this in and solving and simplifying, we end up with the equation of the perpendicular bisector of $AD$:

 $y = \frac{x\sqrt{7}}{7} + \frac{4\sqrt{7}}{7}$.

 
To find the coordinates of $E$, note that $BE$ intersects the perpendicular bisector of $AB$ at $E$ so we just need to solve $\frac{x\sqrt{7}}{3} = \frac{x\sqrt{7}}{7} + \frac{4\sqrt{7}}{7}$ and we get $x = 3$ and so point $E$ is on $(3, \sqrt{7})$.

 
Similarly, to find the coordinates of $F$, note that $CF$ intersects the perpendicular bisector of $AD$ at point $F$ so we just need to solve $-\frac{x\sqrt{7}}{7} + \frac{5\sqrt{7}}{7} = \frac{x\sqrt{7}}{7} + \frac{4\sqrt{7}}{7}$ and we get $x = \frac{1}{2}$. Thus point $F$ is on $(\frac{1}{2}, \frac{9\sqrt{7}}{14})$.

 
Now we have everything we need. We know $A$ is on $(\frac{1}{2}, \frac{3\sqrt{7}}{2})$, $E$ is on $(3, \sqrt{7})$, and $F$ is on $(\frac{1}{2}, \frac{9\sqrt{7}}{14})$.

 We apply Shoelace Theorem to finish the problem off:

 $[AEF] = \frac{1}{2} \mid (\frac{1}{2} \cdot \frac{9\sqrt{7}}{14} + \frac{1}{2} \cdot \sqrt{7} + 3 \cdot \frac{3\sqrt{7}}{2}) - (\frac{3\sqrt{7}}{2} \cdot \frac{1}{2} + \frac{9\sqrt{7}}{14} \cdot 3 + \sqrt{7} \cdot \frac{1}{2}) \mid = \frac{1}{2} \mid \frac{9\sqrt{7}}{28} + \frac{\sqrt{7}}{2} + \frac{9\sqrt{7}}{2} - \frac{3\sqrt{7}}{4} - \frac{27\sqrt{7}}{14} - \frac{\sqrt{7}}{2} \mid = \frac{1}{2} \mid \frac{60\sqrt{7}}{28} \mid = \frac{15\sqrt{7}}{14}$. 

 
Hence, $[AEF] = \frac{15\sqrt{7}}{14}$. Then the answer is just $15 + 7 + 14 = \boxed{036}$.

 




	\item (\textit{AIME I 2020}) Let $P(x)$ be a quadratic polynomial with complex coefficients whose $x^2$ coefficient is $1.$ Suppose the equation $P(P(x))=0$ has four distinct solutions, $x=3,4,a,b.$ Find the sum of all possible values of $(a+b)^2.$


 
 \textbf{Answer: }085



\textbf{Solution 1}

Either $P(3) = P(4)$ or not. We first see that if $P(3) = P(4)$ it's easy to obtain by Vieta's that $(a+b)^2 = 49$. Now, take $P(3) \neq P(4)$ and WLOG $P(3) = P(a), P(4) = P(b)$. Now, consider the parabola formed by the graph of $P$. It has vertex $\frac{3+a}{2}$. Now, say that $P(x) = x^2 - (3+a)x + c$. We note $P(3)P(4) = c = P(3)\left(4 - 4a + \frac{8a - 1}{2}\right) \implies a = \frac{7P(3) + 1}{8}$. Now, we note $P(4) = \frac{7}{2}$ by plugging in again. Now, it's easy to find that $a = -2.5, b = -3.5$, yielding a value of $36$. Finally, we add $49 + 36 = \boxed{085}$. ~awang11, charmander3333

 \textbf{Remark}: We know that $c=\frac{8a-1}{2}$ from $P(3)+P(4)=3+a$.

 





\textbf{Solution 2}

Let the roots of $P(x)$ be $m$ and $n$, then we can write 
 \[P(x)=x^2-(m+n)x+mn\] The fact that $P(P(x))=0$ has solutions $x=3,4,a,b$ implies that some combination of $2$ of these are the solution to $P(x)=m$, and the other $2$ are the solution to $P(x)=n$. It's fairly easy to see there are only $2$ possible such groupings: $P(3)=P(4)=m$ and $P(a)=P(b)=n$, or $P(3)=P(a)=m$ and $P(4)=P(b)=n$ (Note that $a,b$ are interchangeable, and so are $m$ and $n$). We now casework: If $P(3)=P(4)=m$, then 
 \[9-3(m+n)+mn=16-4(m+n)+mn=m \implies m+n=7\]
 \[a^2-a(m+n)+mn=b^2-b(m+n)+mn=n \implies a+b=m+n=7\]so this gives $(a+b)^2=7^2=49$. Next, if $P(3)=P(a)=m$, then 
 \[9-3(m+n)+mn=a^2-a(m+n)+mn=m \implies a+3=m+n\]
 \[16-4(m+n)+mn=b^2-b(m+n)+mn=n \implies b+4=m+n\]Subtracting the first part of the first equation from the first part of the second equation gives 
 \[7-(m+n)=n-m \implies 2n=7 \implies n=\frac{7}{2} \implies m=-3\]Hence, $a+b=2(m+n)-7=2\cdot \frac{1}{2}-7=-6$, and so $(a+b)^2=(-6)^2=36$. Therefore, the solution is $49+36=\boxed{085}$ ~ktong

 





\textbf{Solution 3}

Write $P(x) = x^2+wx+z$. Split the problem into two cases: $P(3)\ne P(4)$ and $P(3) = P(4)$.

 Case 1: We have $P(3) \ne P(4)$. We must have 
 \[w=-P(3)-P(4) = -(9+3w+z)-(16+4w+z) = -25-7w-2z.\]Rearrange and divide through by $8$ to obtain
 \[w = \frac{-25-2z}{8}.\]Now, note that
 \[z = P(3)P(4) = (9+3w+z)(16+4w+z) = \left(9 + 3\cdot \frac{-25-2z}{8} + z\right)\left(16 + 4 \cdot \frac{-25-2z}{8} + z\right) =\]
 \[\left(-\frac{3}{8} + \frac{z}{4}\right)\left(\frac{7}{2}\right) = -\frac{21}{16} + \frac{7z}{8}.\]Now, rearrange to get
 \[\frac{z}{8} = -\frac{21}{16}\]and thus
 \[z = -\frac{21}{2}.\]Substituting this into our equation for $w$ yields $w = -\frac{1}{2}$. Then, it is clear that $P$ does not have a double root at $P(3)$, so we must have $P(a) = P(3)$ and $P(b) = P(4)$ or vice versa. This gives $3+a = \frac{1}{2}$ and $4+b = \frac{1}{2}$ or vice versa, implying that $a+b = 1-3-4 = -6$ and $(a+b)^2 = 36$.

 Case 2: We have $P(3) = P(4)$. Then, we must have $w = -7$. It is clear that $P(a) = P(b)$ (we would otherwise get $P(a)=P(3)=P(4)$ implying $a \in \{3,4\}$ or vice versa), so $a+b=-w=7$ and $(a+b)^2 = 49$.

 Thus, our final answer is $49+36=\boxed{085}$. ~GeronimoStilton

 





\textbf{Solution 4}

Let $P(x) = (x - r)(x - s)$ where $r$ and $s$ are its roots. 

 Note that 
 \[P(r) = P(s) = P(P(3)) = P(P(4)) = P(P(a)) = P(P(b)) = 0.\] Since $P(x)$ is a quadratic, it can have at most two distinct roots, so there must be two pairs of $P(n)$ out of the four $P(3), P(4), P(a),$ and $P(b)$ which are equal to each other. (They have to be pairs because $P(x)$ can only have 2 roots, if we had something like $P(a) = P(3) = P(4) = r$ and $P(b) = s$ then $P(x)$ would have $3$ roots when counting multiplicity which is not possible.)  

 
\textbf{Case 1:} $P(3) = P(4).$

 Then $P(a) = P(b).$ WLOG, let $(3 - r)(3 - s)=(4 - r)(4 - s) = r$ and $(a - r)(a - s)=(b - r)(b - s)=s.$ Solving for $a + b$ gives us $a + b = r + s = 3 + 4 = 7 \implies (a + b)^2 = 49.$

 
\textbf{Case 2:} $P(3) \neq P(4).$ 

 WLOG, let $P(3) = P(a) = r$ and $P(4) = P(b) = s.$ So $(3 - r)(3 - s) = r = (a - r)(a - s)$ and $(4 - r)(4 - s) = s = (b - r)(b - s).$ Subtracting $(4 - r)(4 - s) = s$ from $(3 - r)(3 - s) = r$ gives $-7 + r + s = r -s$, so $s = \frac{7}{2}.$ Plugging this back into $(4-r)(4-s) = s$ gives $r = -3.$ The coefficient of the $x^3$ term in $P(P(x))$ is $-2(r+s)$, so by \href{https://artofproblemsolving.com/wiki/index.php?title=Vieta\%27s\_Formulas}{Vieta's Formulas} the sum of the roots of $P(P(x))$ is $2$ times that of $P(x).$ We know $r + s = \frac{1}{2}$, so $3 + 4 + a + b = \frac{1}{2} \cdot 2 = 1.$ Then $a + b = -6 \implies (a + b)^2 = 36.$ 

 
The requested sum is $36 + 49 =\boxed{85}.$ 

 





\textbf{Solution 5 (Official MAA)}

Note that because $P\big(P(3)\big)=P\big(P(4)\big)= 0$, $P(3)$ and $P(4)$ are roots of $P(x)$. There are two cases.CASE 1: $P(3) = P(4)$. Then $P(x)$ is symmetric about $x=\frac72$; that is to say, $P(r) = P(7-r)$ for all $r$. Thus the remaining two roots must sum to $7$. Indeed, the polynomials $P(x) = \left(x-\frac72\right)^2 + \frac{11}4 \pm i\sqrt3$ satisfy the conditions.CASE 2: $P(3)\neq P(4)$. Then $P(3)$ and $P(4)$ are the two distinct roots of $P(x)$, so
 \[P(x) = \big(x-P(3)\big)\big(x-P(4)\big)\]for all $x$. Note that any solution to $P\big(P(x)\big) = 0$ must satisfy either $P(x) = P(3)$ or $P(x) = P(4)$. Because $P(x)$ is quadratic, the polynomials $P(x) - P(3)$ and $P(x) - P(4)$ each have the same sum of roots as the polynomial $P(x)$, which is $P(3) + P(4)$. Thus the answer in this case is $2\big(P(3) + P(4)\big)-7$, and so it suffices to compute the value of $P(3)+P(4)$.

 Let $P(3)=u$ and $P(4) = v$. Substituting $x=3$ and $x=4$ into the above quadratic polynomial yields the system of equations
 \begin{align*} u &= (3-u)(3-v) = 9 - 3u - 3v + uv\\ v &= (4-u)(4-v) = 16 - 4u - 4v + uv. \end{align*}
Subtracting the first equation from the second gives $v - u = 7 - u - v$, yielding $v = \frac72.$ Substituting this value into the second equation gives
 \[\dfrac72 = \left(4 - u\right)\left(4 - \dfrac72\right),\]yielding $u = -3.$ The sum of the two solutions is $2\left(\frac72-3\right)-7 = -6$. In this case, $P(x)= (x+3)\left(x-\frac72\right)$.

 The requested sum of squares is $7^2+(-6)^2 = \boxed{085}$.

 





\textbf{Solution 6}

Let $P(x) = (x-c)^2 - d$ for some $c$, $d$.

 Then, we can write $P(P(x)) = ((x-c)^2 - d - c)^2 - d$. Setting the expression equal to $0$ and solving for $x$ gives:

 $x = \pm \sqrt{ \pm \sqrt{d} + d + c} + c$

 Therefore, we have that $x$ takes on the four values $\sqrt{\sqrt{d} + d + c} + c$, $-\sqrt{\sqrt{d} + d + c} + c$, $\sqrt{-\sqrt{d} + d + c} + c$, and $-\sqrt{-\sqrt{d} + d + c} + c$. Two of these values are $3$ and $4$, and the other two are $a$ and $b$.

 We can split these four values into two "groups" based on the radicand in the expression - for example, the first group consists of the first two values listed above, and the second group consists of the other two values.

 $\textbf{Case 1}$: Both the 3 and 4 values are from the same group.

 In this case, the $a$ and $b$ values are both from the other group. The sum of this is just $2c$ because the radical cancels out. Because of this, we can see that $c$ is just the average of $3$ and $4$, so we have $2c = 3 + 4 = 7$, so $(a+b)^2 = 7^2 = 49$.

 $\textbf{Case 2}$: The 3 and 4 values come from different groups.

 It is easy to see that all possibilities in this case are basically symmetric and yield the same value for $(a+b)^2$. Without loss of generality, assume that $\sqrt{\sqrt{d} + d + c} + c = 4$ and $\sqrt{-\sqrt{d} + d + c} + c = 3$. Note that we can't switch the values of these two expressions since the first one is guaranteed to be larger.

 We can write $\sqrt{\sqrt{d} + d + c} + c = 1 + \sqrt{-\sqrt{d} + d + c} + c$.

 Moving most terms to the left side and simplifying gives $\sqrt{\sqrt{d} + d + c} - \sqrt{-\sqrt{d} + d + c} = 1$.

 We can square both sides and simplify:

 $\sqrt{d} + d + c - \sqrt{d} + d + c - 2\sqrt{(d + c + \sqrt{d})(d + c - \sqrt{d})} = 1$

 $2d + 2c - 2\sqrt{(d + c + \sqrt{d})(d + c - \sqrt{d})} = 1$

 $\sqrt{(d + c + \sqrt{d})(d + c - \sqrt{d})} = (d+c) - \frac{1}{2}$

 $\sqrt{(d+c)^2 - (\sqrt{d})^2} = (d+c) - \frac{1}{2}$

 $\sqrt{d^2 + 2dc + c^2 - d} = (d+c) - \frac{1}{2}$

 Squaring both sides again gives the following:

 $d^2 + 2dc + c^2 - d = d^2 + 2dc + c^2 - d - c + \frac{1}{4}$

 Nearly all terms cancel out, yielding $c = \frac{1}{4}$.

 By substituting this back in, we obtain $\sqrt{\sqrt{d} + d + c} = \frac{15}{4}$ and $\sqrt{-\sqrt{d} + d + c} = \frac{11}{4}$.

 The sum of $a$ and $b$ is equal to $-\sqrt{\sqrt{d} + d + c} - \sqrt{-\sqrt{d} + d + c} + 2c = -\frac{15}{4} - \frac{11}{4} + \frac{1}{2} = -6$, so $(a+b)^2 = 36$.

 Adding up both values gives $49 + 36 = \boxed{085}$ as our final answer.

 


 


 




	\item (\textit{AIME I 2020}) Let $\triangle ABC$ be an acute triangle with circumcircle $\omega,$ and let $H$ be the intersection of the altitudes of $\triangle ABC.$ Suppose the tangent to the circumcircle of $\triangle HBC$ at $H$ intersects $\omega$ at points $X$ and $Y$ with $HA=3,HX=2,$ and $HY=6.$ The area of $\triangle ABC$ can be written in the form $m\sqrt{n},$ where $m$ and $n$ are positive integers, and $n$ is not divisible by the square of any prime. Find $m+n.$


 
 \textbf{Answer: }058



\textbf{Solution 1}

The following is a power of a point solution to this menace of a problem:

 
                    \begin{figure}[H]
                        \centering
                        \includegraphics[width=0.5\textwidth]{Solution/AIME_2020_I_15_Solution_Figure_1.png}
                    \end{figure}
                    Let the points in the diagram be what they appear to be. Note that $3HX = HY$ is not insignificant; from here, we set $XH = HE = \frac{1}{2} EY = HL = 2$ by PoP and trivial construction. Now, $D$ is the reflection of $A$ over $H$. Note $AO \perp XY$, and therefore by Pythagorean theorem we have $AE = XD = \sqrt{5}$. Consider $HD = 3$. We have that $\triangle HXD \cong HLK$, and therefore we are ready to PoP with respect to $(BHC)$. Setting $BL = x, LC = y$, we obtain $xy = 10$ by PoP on $(ABC)$, and furthermore, we have $KH^2 = 9 = (KL - x)(KL + y) = (\sqrt{5} - x)(\sqrt{5} + y)$. Now, we get $4 = \sqrt{5}(y - x) - xy$, and from $xy = 10$ we take 
 \[\frac{14}{\sqrt{5}} = y - x.\] However, squaring and manipulating with $xy = 10$ yields that $(x + y)^2 = \frac{396}{5}$ and from here, since $AL = 5$ we get the area to be $3\sqrt{55} \implies \boxed{058}$. ~awang11's sol

 





\textbf{Solution 1a}

As in the diagram, let ray $AH$ extended hits BC at L and the circumcircle at say $P$. By power of the point at H, we have $HX \cdot HY = AH \cdot HP$. The three values we are given tells us that $HP=\frac{2\cdot 6}{3}=4$. L is the midpoint of $HP$(see here: \href{https://www.cut-the-knot.org/Curriculum/Geometry/AltitudeAndCircumcircle.shtml}{https://www.cut-the-knot.org/Curriculum/Geometry/AltitudeAndCircumcircle.shtml} ), so $HL=LP=2$.

 As in the diagram provided, let K be the intersection of $BC$ and $XY$. By power of a point on the circumcircle of triangle $HBC$, $KH^{2}=KB \cdot KC$. By power of a point on the circumcircle of triangle $ABC$, $KB \cdot KC=KX \cdot KY$, thus $KH^{2}=(KH-2)(KH+6)$. Solving gives $4KH=12$ or $KH=3$.

 By the Pythagorean Theorem on triangle $HKL$, $KL=\sqrt{5}$. Now continue with solution 1.

 





\textbf{Solution 2}


                    \begin{figure}[H]
                        \centering
                        \includegraphics[width=0.5\textwidth]{Solution/AIME_2020_I_15_Solution_Figure_1.png}
                    \end{figure}
                    Diagram not to scale.

 
We first observe that $H'$, the image of the reflection of $H$ over line $BC$, lies on circle $O$.  This is because $\angle HBC = 90 - \angle C = \angle H'AC = \angle H'BC$.  This is a well known lemma.  The result of this observation is that circle $O'$, the circumcircle of $\triangle BHC$ is the image of circle $O$ over line $BC$, which in turn implies that $\overline{AH} = \overline{OO'}$ and thus $AHO'O$ is a parallelogram.  That $AHO'O$ is a parallelogram implies that the radius $AO$ is perpendicular to $\overline{XY}$, and thus divides segment $\overline{XY}$ in two equal pieces, $\overline{XD}$ and $\overline{DY}$, of length $4$.

 
Using Power of a Point,
 \[\overline{AH} \cdot \overline{HH'} = \overline{XH} \cdot \overline{HY} \Longrightarrow 3 \cdot \overline{HH'} = 2 \cdot 6 \Longrightarrow \overline{HH'} = 4\]This means that $\overline{HL} = \frac12 \cdot 4 = 2$ and $\overline{AL} = 2 + 3 = 5$, where $L$ is the foot of the altitude from $A$ onto $BC$.  All that remains to be found is the length of segment $\overline{BC}$.

 Looking at right triangle $\triangle AHD$, we find that
 \[\overline{AD} = \sqrt{\overline{AH}^2 - \overline{HD}^2} = \sqrt{3^2 - 2^2} = \sqrt{5}\]Looking at right triangle $\triangle ODY$, we get the equation
 \[\overline{OY}^2 - \overline{DY}^2 = \overline{OD}^2 = \left(\overline{AO} - \overline{AD}\right)^2\]Plugging in known values, and letting $R$ be the radius of the circle, we find that
 \[R^2 - 16 = (R - \sqrt{5})^2 = R^2 - 2\sqrt5 R + 5 \Longrightarrow R = \frac{21\sqrt5}{10}\]

 Recall that $AHO'O$ is a parallelogram, so $\overline{AH} = \overline{OO'} = 3$.  So, $\overline{OM} = \frac32$, where $M$ is the midpoint of $\overline{BC}$.  This means that
 \[\overline{BC} = 2\overline{BM} = 2\sqrt{R^2 - \left(\frac32\right)^2} = 2\sqrt{\frac{441}{20} - \frac{9}{4}} = \frac{6\sqrt{55}}{5}\]

 Thus, the area of triangle $\triangle ABC$ is
 \[\frac{\overline{AL} \cdot \overline{BC}}{2} = \frac{5 \cdot \frac{6\sqrt{55}}{5}}{2} = {3\sqrt{55}}\]The answer is $3 + 55 = \boxed{058}$.

 





\textbf{Solution 3 (Official MAA 1)}

Extend $\overline{AH}$ to intersect $\omega$ again at $P$. The Power of a Point Theorem yields $HP = \frac{HX \cdot HY}{HA} = 4$. Because $\angle CAP=\angle CBP$, and $\angle CAP$ and $\angle CBH$ are both complements to $\angle C$, it follows that $\angle CBP = \angle CBH$, implying that $\overline{BC}$ bisects $\overline{HP}$, so the length of the altitude from $A$ to $\overline{BC}$ is $h_a = AH + \frac12 HP = 5$.

 Let the circumcircle of $\triangle BCH$ be $\omega'$. Because $\triangle BCH \cong \triangle BCP$, the two triangles must have the same circumradius. Because the circumcircle of $\triangle BCP$ is $\omega$, the circles $\omega$ and $\omega'$ have the same radius $R$. Denote the centers of $\omega$ and $\omega'$ by $O$ and $O'$, respectively, and let $M$ be the midpoint of $\overline{XY}$. Note that trapezoid $HMOO'$ has $\angle H = \angle M = 90^\circ$. Also $HM = XM - XH = \frac12\cdot XY - HX = 2$ and $HO' = R$. Because $\omega$ is a translation of $\omega'$ in the direction of $\overline{AH}$, it follows that $OO' = AH = 3$. Finally, the Pythagorean Theorem applied to $\triangle XMO$ yields $MO = \sqrt{R^2-16}$. Let $T$ be the projection of $O$ onto $\overline{HO'}$. Then $TO' = R-MO$, so the Pythagorean Theorem applied to $\triangle TOO'$ yields
 \[R - \sqrt{R^2-16} = \sqrt{3^2 - 2^2} = \sqrt{5}.\]Solving for $R$ gives $R = \frac{21}{2\sqrt5}$. It follows from properties of the orthocenter $H$ that
 \[\cos\angle A = \dfrac{AH}{2R} = \dfrac{\sqrt5}{7},\]so
 \[\sin\angle A = \sqrt{1 - \cos^2\angle A} = \dfrac{2\sqrt{11}}{7}.\]Therefore by the Extended Law of Sines
 \[a = BC = 2R \sin\angle A = \dfrac{6\sqrt{11}}{\sqrt5},\]so
 \[[\triangle ABC] = \frac12 a h_a = \frac12 \cdot \frac{6\sqrt{11}}{\sqrt{5}} \cdot 5 = 3\sqrt{55}.\]The requested sum is $3+55 = 58$.

 
                    \begin{figure}[H]
                        \centering
                        \includegraphics[width=0.5\textwidth]{Solution/AIME_2020_I_15_Solution_Figure_1.png}
                    \end{figure}
                    

 





\textbf{Solution 3a (slightly modified Solution 3)}

Note that $\overline{BC}$ bisects $\overline{OO'}$. Using the same method from Solution 3 we find $R = \frac{21}{2\sqrt{5}}$. Let the midpoint of $\overline{BC}$ be $N$, then by the Pythagorean Theorem we have $BN^2 = OB^2 - ON^2 = R^2 - \frac{3}{2}^2$, so $BN = \frac{3\sqrt{11}}{\sqrt{5}}$. Since $h_a = 5$ we have that the area of ABC is $3\sqrt{55}$ so the answer is $3 + 55 = \boxed{58}$

 





\textbf{Solution 4 (Official MAA 2)}

Let $D$ be the intersection point of line $AH$ and $\overline{BC}$, noting that $\overline{AD}\perp\overline{BC}$. Because the area of $\triangle ABC$ is $\frac12\cdot AD\cdot BC$, it suffices to compute $AD$ and $BC$ separately. As in the previous solution, $AD = 5$. The value of $BC$ can be found using the following lemma.

 Lemma: Triangle $AXY$ is isosceles with base $\overline{XY}$.

 Proof: Because the circumcircle of $\triangle BCH$, $\omega'$, and $\omega$ have the same radius, there exists a translation $\Phi$ sending the former to the latter. Because $\overline{AH}$ is parallel to the line connecting the centers of the two circles, $\Phi$ must send $H$ to $A$, meaning $\Phi$ also sends $\overline{XY}$ to the tangent to $\omega$ at $A$. But this means that this tangent is parallel to $\overline{XY}$, which implies the conclusion.

 Applying Stewart's Theorem to $\triangle AXY$ yields
 \[AX^2 = AH^2 + HX\cdot HY = 3^2 + 2\cdot 6 = 21,\]implying $AX = AY= \sqrt{21}.$

 By the Law of Cosines
 \[\cos \angle XAY = \frac{21 + 21 - 64}{2 \cdot 21} = -\frac{11}{21},\]so
 \[\sin \angle XAY = \dfrac{4\sqrt{20}}{21}.\]Let $R$ be the radius of $\omega$. By the Extended Law of Sines
 \[R = \frac{XY}{2 \cdot \sin \angle XAY} = \dfrac{21}{\sqrt{20}}.\]Then the solution proceeds as in Solution 3.

 


 





\textbf{Solution 5 (Official MAA 3)}

Define points $D$ and $P$ as above, and note that $AD=5$ and $DH=PD=AD - AH = 2$. Let the circumcircle of $\triangle BCH$ be $\omega'$.

 Extend $\overline{XY}$ past $X$ until it intersects line $BC$ at point $Z$. Because line $BC$ is a radical axis of $\omega$ and $\omega'$, it follows from the Power of a Point Theorem that
 \[ZX \cdot ZY = ZX \cdot(ZX + 8) = ZH^2 = (ZX+2)^2,\]from which $ZX=1$. By Pythagorean Theorem $ZD=\sqrt{ZH^2 - DH^2} = \sqrt{5}$.

 


 Let $m=CD$ and $n=BD$. By the Power of a Point Theorem
 \[mn= AD\cdot PD = 10.\]On the other hand,
 \[ZH^2 = 9 = ZB \cdot ZC = (\sqrt{5} - n)(\sqrt{5} + m) = 5 + \sqrt{5}(m-n) - mn,\]from which $m-n = \frac{14}{\sqrt{5}}$. Therefore

 
 \[(m+n)^2 = (m-n)^2 + 4mn = \frac{196}{5} + 40 = \frac{396}{5}.\]Thus $BC = m+n=\sqrt{\frac{396}{5}} = 6\sqrt{\frac{11}{5}}$. Therefore$[\triangle ABC] = \frac{1}{2}BC \cdot AD = 3\sqrt{55},$ as above.

 





\textbf{Solution 6}


                    \begin{figure}[H]
                        \centering
                        \includegraphics[width=0.5\textwidth]{Solution/AIME_2020_I_15_Solution_Figure_1.png}
                    \end{figure}
                    Let $O$ be circumcenter of $ABC,$ let $R$ be circumradius of $ABC,$ let $\omega'$ be the image of circle $\omega$ over line $BC$ (the circumcircle of $HBC$). 

 Let $P$ be the image of the reflection of $H$ over line $BC, P$ lies on circle $\omega.$ Let $M$ be the midpoint of $XY.$Then $P$ lies on $\omega, OA = O'H, OA || O'H.$

 (see here: \href{https://brilliant.org/wiki/triangles-orthocenter/}{https://brilliant.org/wiki/triangles-orthocenter/}or \href{https://en.wikipedia.org/wiki/Altitude\_(triangle)}{https://en.wikipedia.org/wiki/Altitude\_(triangle)} Russian)

 $P$ lies on $\omega \implies OA = OP = R,$ 

 $OA = O'H, OA || O'H \implies M$ lies on $OA.$

 We use properties of crossing chords and get 
 \[AH \cdot HP = XH \cdot HY = 2 \cdot 6 \implies HP = 4, AP = AH + HP = 7.\]We use properties of radius perpendicular chord and get
 \[MH = \frac{XH + HY}{2} – HY = 2.\]We find
 \[\sin OAH =\frac{MH}{AH} = \frac{2}{3}  \implies  \cos OAH = \frac{\sqrt{5}}{3}.\]We use properties of isosceles $\triangle OAP$ and find $\hspace{5mm}R = \frac{AP}{2\cos OAP} = \frac{7}{2\frac {\sqrt{5}}{3}} = \frac{21}{2\sqrt{5}}.$

 We use $OM' = \frac{AH}{2} = \frac {3}{2}$ and find $\hspace{25mm} \frac{BC}{2} = \sqrt{R^2 – OM'^2} = 3 \sqrt {\frac {11}{5}}.$

 The area of $ABC$
 \[[ABC]=\frac{BC}{2} \cdot (AH + HD) = 3\cdot \sqrt{55}  \implies 3+55 = \boldsymbol{\boxed{058}}.\]\textbf{vladimir.shelomovskii@gmail.com, vvsss}

 \textbf{Two Video Solutions by MOP 2024} 
                \begin{itemize}
                    \item \href{https://youtube.com/watch?v=qgM-GbAWXCc}{https://youtube.com/watch?v=qgM-GbAWXCc} (Power of a Point and Symmetry)\item \href{https://youtube.com/watch?v=XOUssYh14aw}{https://youtube.com/watch?v=XOUssYh14aw} (Geometric Transformations, Angle Chasing, Trig)
                \end{itemize}


                




\end{enumerate}